\documentclass[11pt,a4paper]{article}

\usepackage[utf8]{inputenc}
\usepackage[T1]{fontenc}
\usepackage{amsmath,amssymb,amsthm,mathtools}
\usepackage{graphicx}
\graphicspath{{../../validation-federated/}{./}}
\usepackage{booktabs}
\usepackage{array}
\usepackage{enumitem}
\usepackage{float}
\usepackage{listings}
\usepackage{xcolor}
\usepackage{siunitx}
\usepackage{microtype}
\usepackage{lmodern}
\usepackage[margin=2.4cm]{geometry}
\usepackage[hidelinks]{hyperref}
\usepackage{cleveref}
\usepackage{natbib}
\usepackage{stmaryrd}

\theoremstyle{plain}
\newtheorem{theorem}{Theorem}[section]
\newtheorem{lemma}[theorem]{Lemma}
\newtheorem{proposition}[theorem]{Proposition}
\newtheorem{corollary}[theorem]{Corollary}
\theoremstyle{definition}
\newtheorem{definition}[theorem]{Definition}
\newtheorem{assumption}[theorem]{Assumption}
\newtheorem{construction}[theorem]{Construction}
\newtheorem{principle}[theorem]{Principle}
\theoremstyle{remark}
\newtheorem{remark}[theorem]{Remark}
\newtheorem{example}[theorem]{Example}
\newtheorem{observation}[theorem]{Observation}

%--- notation -----------------------------------------------------------
\newcommand{\Plan}{\mathcal{P}}
\newcommand{\Step}{\sigma}
\newcommand{\Src}{\mathsf{S}}
\newcommand{\Capset}{\mathrm{cap}}
\newcommand{\Req}{\mathrm{req}}
\newcommand{\Low}{\mathrm{low}}
\newcommand{\Ext}{\mathrm{ext}}
\newcommand{\Res}{\mathcal{R}}
\newcommand{\Bind}{\Gamma}
\newcommand{\Verd}{\mathrm{v}}
\newcommand{\vans}{\textsf{answered}}
\newcommand{\vempty}{\textsf{empty}}
\newcommand{\vsurf}{\textsf{unsupported}}
\newcommand{\vtime}{\textsf{exhausted}}
\newcommand{\vref}{\textsf{refused}}
\newcommand{\vstarve}{\textsf{starved}}
\newcommand{\bmdl}{\textsf{blocked-by-model}}
\newcommand{\bcorp}{\textsf{blocked-by-corpus}}
\newcommand{\beng}{\textsf{blocked-by-engine}}
\newcommand{\bbud}{\textsf{blocked-by-budget}}
\newcommand{\Univ}{\mathcal{U}}
\newcommand{\Nspace}{\mathcal{N}}
\newcommand{\idm}{\mathrm{id}}
\newcommand{\Rset}{\mathbb{R}}
\newcommand{\Nset}{\mathbb{N}}
\newcommand{\Bud}{B}
\newcommand{\cost}{c}
\newcommand{\HFQ}{\textsc{hfq}}
\newcommand{\Feat}{\mathcal{F}}
\newcommand{\Zset}{\mathbb{Z}}

\lstdefinestyle{hfq}{
  basicstyle=\ttfamily\footnotesize,
  keywordstyle=\bfseries,
  commentstyle=\itshape\color{gray!70!black},
  stringstyle=\color{black},
  showstringspaces=false,
  breaklines=true,
  frame=single,
  framerule=0.3pt,
  rulecolor=\color{gray!60},
  numbers=left,
  numberstyle=\tiny\color{gray},
  xleftmargin=2.2em,
  morekeywords={let,from,ask,where,filter,join,map,via,keep,drop,union,
    intersect,minus,require,budget,within,on,fail,else,emit,assert,using,
    plan,source,expect,cost,limit,order,by,as,in,not,and,or,exists,
    with,when,unresolved,partial}
}
\lstset{style=hfq}

\title{\textbf{A Plan Language for Federated Biological Retrieval}\\[0.4em]
\large Capability-indexed lowering, verdict-carrying steps,\\
and the arithmetic of partial identifier translation}

\author{%
Kundai Farai Sachikonye\\[0.3em]
\normalsize Experimental Bioinformatics, TUM School of Life Sciences\\
\normalsize Maximus-von-Imhof-Forum 3, 85354 Freising, Germany\\
\normalsize \texttt{kundai.sachikonye@bitspark.com}}

\date{\today}

\begin{document}
\maketitle

\begin{abstract}
A biological question that spans more than one public database is today
answered by a script: a host-language program that interpolates identifiers
into query strings, sends them, parses what returns, and interpolates again.
The host language knows nothing about the data and the query language knows
nothing about control flow, and every failure mode of federated retrieval lives
in the seam between them. We give a language that closes the seam by refusing
to be a query language at all.

The design rests on one observation. The databases a biologist must cross ---
a reaction knowledge base with a SPARQL endpoint, a protein resource exposing
both SPARQL and a REST interface, a pathway resource exposing only flat-file
REST, an ontology available as a downloadable artefact --- do not share a query
model, and no amount of engineering will give them one. They do share a
\emph{result} model: finite sets of identifiers carrying attributes. A language
whose terms denote result sets and whose operators denote instructions on what
to do with a result is therefore portable across backends in a way that a
language whose terms denote graph patterns can never be. We call the resulting
object a \emph{plan language}; queries are its leaves and are never written by
the user.

We develop the language in four layers. \Cref{sec:model} fixes a source model
in which a backend is a triple of an endpoint, a capability set, and an
extraction function, and proves that lowering a plan step is total exactly on
the steps whose required capabilities are contained in the source's declared
set (\Cref{thm:lowering}). \Cref{sec:ir} gives the intermediate representation
and shows that a plan admits a static capability check that decides
compilability before any request is issued (\Cref{thm:static}), which converts
a class of runtime failures into compile-time diagnostics.
\Cref{sec:verdicts} replaces the boolean answer of conventional query
interfaces with a six-valued verdict and proves that the six are pairwise
distinguishable by a plan executor while any two of them are conflated by an
interface returning rows alone (\Cref{thm:six}, \Cref{cor:onebit}). The
verdict $\vstarve$ --- a step that failed because an \emph{earlier} step
under-retrieved --- is reachable only in the federated setting, and we show
that it is the characteristic failure of federation rather than an exotic case
(\Cref{prop:starve-reachable}).
\Cref{sec:translation} treats identifier translation, which is where the real
losses occur: cross-namespace maps are partial, non-injective and
non-confluent, and we give exact composition laws for retention
(\Cref{thm:retention}), a proof that translation loss is not recoverable by
reordering (\Cref{thm:reorder}) and a proof that two plans differing only in
translation route may return different sets without either being wrong
(\Cref{thm:route-extent}).

Three further results concern execution. \Cref{thm:allocation} solves budget
allocation across heterogeneous sources by a water-filling argument with a
single shadow price, and \Cref{prop:necessary-not-sufficient} shows that a
budget exceeding the sum of minimum step costs is necessary but not sufficient
for a plan to complete. \Cref{thm:interpolation} shows that structural
interpolation --- lowering a bound result set into a target language through
the language's own value-binding construct rather than through string
concatenation --- eliminates by construction a family of semantic defects that
textual interpolation admits, and we exhibit a documented instance of that
family in which two spellings of the same intent against the same endpoint
returned counts differing by more than two orders of magnitude.
\Cref{thm:idempotent} characterises when a plan may be safely re-executed.

We specify an executable prototype, report its behaviour on offline fixtures
across four source kinds, and state plainly what has not been established: the
prototype has not been run against any live public endpoint, no claim is made
about answer correctness against biology, and the capability declarations are
asserted by the source adapter rather than verified against the service.
\end{abstract}

\tableofcontents

\clearpage

%==============================================================================
\part{The Setting}
%==============================================================================

\section{Introduction}
\label{sec:intro}

\subsection{The script that everybody writes}

Consider a question of the kind that arises daily in metabolic research: which
enzymes catalyse reactions consuming any member of a given chemical class, and
of those, which occur in pathways that also contain a second class. No single
public resource answers it. The reaction membership lives in one database, the
enzyme--reaction association in a second, the pathway membership in a third,
and the chemical class hierarchy in a fourth.

The universal solution is a script. In outline:

\begin{lstlisting}[caption={The script everybody writes. Not a strawman.},label=lst:script]
compounds = sparql(ENDPOINT_A, TEMPLATE_A % class_iri)
ids       = [row["c"] for row in compounds]
enzymes   = []
for chunk in batches(ids, 50):
    q = TEMPLATE_B % " ".join("<%s>" % i for i in chunk)
    enzymes += sparql(ENDPOINT_A, q)
kegg_ids  = [xref(e) for e in enzymes]
pathways  = [rest_get(ENDPOINT_C, "/link/pathway/%s" % k) for k in kegg_ids]
\end{lstlisting}

Every practitioner recognises \cref{lst:script}, and every practitioner
dislikes it. The dislike is usually expressed aesthetically. It is not an
aesthetic problem. \Cref{lst:script} exhibits, in eight lines, five distinct
and independently serious defects.

\begin{enumerate}[label=(D\arabic*),leftmargin=3em]
\item \textbf{The interpolation is textual.} Line~5 constructs a query by
string concatenation. Whether the resulting string means what the author
intended depends on the target language's parse of a string the author never
sees in full. \Cref{sec:interp} exhibits a documented case in which two
spellings of one intent, differing in a single line, returned $2$ and $397$
against the same service on the same day.
\item \textbf{Failure is one bit wide.} If \texttt{compounds} comes back empty,
the script cannot distinguish: the class has no members; the endpoint timed out
and returned an empty result rather than an error; the query used a construct
the endpoint does not support and the endpoint degraded silently; the class
identifier was misspelled. Each demands a different repair, and the script sees
one empty list.
\item \textbf{The failure of step $n$ is attributed to step $n$.} If
\texttt{pathways} is empty, the proximate cause is very often that
\texttt{kegg\_ids} was short, because \texttt{xref} is a partial function and
silently dropped two-fifths of its input. The script reports a pathway problem
and the analyst debugs the wrong stage.
\item \textbf{The budget is unstated and unallocated.} There is a rate limit on
\texttt{ENDPOINT\_C}, a timeout on \texttt{ENDPOINT\_A}, and no representation
of either anywhere in the program. The script discovers them by failing.
\item \textbf{The control flow is invisible to the data layer.} The loop over
\texttt{batches} exists because a query language has no way to say ``and
then''. The batching constant $50$ is a magic number encoding an unstated
belief about request-size limits at a service the script does not name.
\end{enumerate}

These are not five instances of sloppiness. They are five consequences of one
structural decision --- that control flow lives in a host language and data
access lives in a query language, and neither can see the other.

\subsection{Why the obvious repair fails}

The obvious repair is to make the query language expressive enough to hold the
whole computation: one large query with unions, optional patterns, subselects
and federated service clauses. This is what the SPARQL~1.1 federation extension
was designed for \citep{Prudhommeaux2013Federated,Harris2013}, and it fails for
two reasons that are not implementation defects.

First, it is frequently not available. Of the four sources in the motivating
question, only two speak SPARQL at all. A pathway resource exposing a flat-file
REST interface \citep{Kanehisa2000,Kanehisa2023} cannot be reached by a service
clause, because there is no SPARQL protocol endpoint to address. The federation
extension federates SPARQL endpoints, and the biological data landscape is not
made of SPARQL endpoints.

Second, where it is available, it concentrates the defects rather than removing
them. A single query containing every stage of the computation is a single
query whose failure is still one bit wide (D2), whose intermediate results are
still invisible (D3), and which is now large enough that the interpolation
defect (D1) becomes very hard to see --- as \cref{sec:interp} demonstrates on a
query of eight lines.

A third repair, the workflow system \citep{Wolstencroft2013,Koster2012,Berthold2009},
supplies exactly the control flow (D5) that the query language lacks. It leaves
(D1)--(D4) untouched, because a workflow node that issues a query issues a
string that its author assembled, and the workflow engine does not know what
the string means.

\subsection{What we build instead}

We separate the two things the script conflates. A \emph{plan} is a sequence of
\emph{steps}; a step names a source, an abstract pattern to send it, and a
binding to receive the result; the plan language supplies the control flow and
the source adapters supply the concrete requests. The user writes the plan.
Nobody writes a query.

\begin{principle}[Queries are leaves]
\label{prin:leaves}
No construct of the plan language denotes a graph pattern, a triple, a
relational clause, or a URL. Every such object is produced by \emph{lowering},
from an abstract step together with a source declaration, and is not
addressable by the plan author.
\end{principle}

\Cref{prin:leaves} is what makes the language portable across backends that
share no query model. Its cost is stated immediately: a plan cannot express a
request that its source adapters cannot generate, and \cref{sec:limits} records
what that excludes.

\begin{principle}[A step returns a verdict, not rows]
\label{prin:verdict}
The value of a step is a pair (verdict, payload). The payload is a result set
only when the verdict is $\vans$. Every other verdict carries a diagnosis and
no result set.
\end{principle}

\Cref{prin:verdict} is what makes the language diagnostic. Its content is
\cref{sec:verdicts}, and the reason it is not a mere convention is
\cref{cor:onebit}: an interface that returns rows, and signals failure by
returning zero of them, cannot represent the distinction at all, whatever
discipline is imposed on its callers.

\begin{principle}[Refusal is a first-class outcome]
\label{prin:refusal}
A plan whose declared requirements exceed what its sources can supply, or whose
budget falls below the minimum the plan needs, terminates without issuing any
request and reports what was missing.
\end{principle}

A system with no refusals declines nothing, and therefore claims everything it
is asked. \Cref{prin:refusal} is the operational form of that observation, and
\cref{thm:static} is the result that makes it enforceable statically rather
than by convention.

\subsection{Contributions and non-contributions}

We claim: a source model under which heterogeneous backends compose
(\cref{sec:model}); a capability calculus deciding compilability before
execution (\cref{thm:static}); a six-valued verdict algebra with a proof that
the conventional interface conflates its members (\cref{thm:six},
\cref{cor:onebit}); exact laws for retention under partial identifier
translation (\cref{thm:retention}) with a proof that reordering does not
recover loss (\cref{thm:reorder}) and that two translation routes may disagree
without either being wrong (\cref{thm:route-extent}); a budget allocation with
a single shadow price (\cref{thm:allocation}) and a proof that sufficient total
budget is not sufficient (\cref{prop:necessary-not-sufficient}); a proof that
structural interpolation eliminates a defect family that textual interpolation
admits (\cref{thm:interpolation}); and a characterisation of safe re-execution
(\cref{thm:idempotent}).

We do not claim: that plans written in this language return biologically
correct answers; that the capability declarations of \cref{sec:model} are
verified rather than asserted; that the retention laws of
\cref{sec:translation} have been measured against real cross-reference tables;
or that the prototype of \cref{sec:proto} has been exercised against any live
service. \Cref{sec:limits} enumerates these and others, and
\cref{rem:honesty-assumption} states the single assumption on which the largest
part of the development rests.

\subsection{Conventions}

$\Univ$ denotes a countable universe of identifiers. A \emph{namespace} is a
subset of $\Univ$; namespaces are pairwise disjoint, which is a modelling
decision and not a fact about the world (\cref{rem:disjoint}). Result sets are
finite. We write $f : X \rightharpoonup Y$ for a partial function and
$\operatorname{dom} f$ for its domain of definition. All complexity statements
count \emph{requests} unless stated otherwise, since request count is what a
public service rate-limits and what a plan can control; wall-clock time is not
modelled, and \cref{sec:limits} records the consequence.

\section{Related work and where this sits}
\label{sec:related}

Federated query processing over distributed data has a substantial literature
concerned principally with \emph{source selection} and \emph{join ordering}:
given a query and a set of endpoints, decide which endpoints can contribute and
in what order to contact them \citep{Schwarte2011,Gorlitz2011,Acosta2011,Saleem2016,Rakhmawati2013}.
That work assumes a single query language across sources and optimises within
it. Our setting is complementary: the sources do not share a language, the
decomposition is written by the user rather than inferred, and the object of
study is what a step reports when it fails rather than how fast it succeeds.
The two concerns are orthogonal, and a planner of the kind those papers
describe could in principle choose among the routes that \cref{sec:translation}
shows to be inequivalent --- but it would need a cost model over partial maps
that the literature does not supply.

Mediator and wrapper architectures \citep{Wiederhold1992,GarciaMolina1997,Haas1997}
established the pattern of a per-source adapter presenting a common interface,
and capability-based query planning \citep{Levy1996,Papakonstantinou1995,Florescu1999,Yerneni1999}
formalised sources that can answer only some queries, binding-pattern
restrictions being the classical instance. Our capability sets
(\cref{def:capability}) are in that lineage, and \cref{thm:static} is a
static-check analogue of their planning-time feasibility tests. The difference
lies in what happens on failure: those systems treat an uncoverable query as a
planning failure, and we treat it as a typed verdict that the plan can catch
and route around, which is the difference between a compiler error and an
exception.

Scientific workflow systems \citep{Wolstencroft2013,Berthold2009,Koster2012,DiTommaso2017}
provide the control flow we require, and are widely used for exactly the
motivating question. They are host languages with library calls, so the requests
they issue are strings assembled by the workflow author, and defects
(D1)--(D4) survive unchanged. A workflow system sequences steps; it does not
know what a step means, and so cannot report why one came back empty.

Identifier translation across biological namespaces has its own literature and
its own well-documented pathologies \citep{Chen2005,VanIersel2010,Wimalaratne2018,Juty2012,Ison2016}.
The standard framing is pairwise coverage: what fraction of namespace $A$ maps
into namespace $B$. \Cref{sec:translation} takes the \emph{composition} of
several such maps as the object, and derives what survives a chain --- a
quantity a plan actually depends on and which is not determined by the pairwise
coverages alone (\cref{prop:pairwise-insufficient}).

Finally, the observation that an empty result and a failure are different, and
that conflating them corrupts downstream reasoning, is old in database theory:
it is the null-value problem \citep{Codd1979,Imielinski1984,Libkin2016} in a new
setting. \Cref{sec:verdicts} is a six-valued refinement specific to federation,
and \cref{prop:starve-reachable} identifies a value with no analogue in the
single-database case.

%==============================================================================
\part{The Language}
%==============================================================================

\section{The source model}
\label{sec:model}

\subsection{Sources differ in kind, not dialect}

The design turns on a fact about the data landscape that is easy to state and
consequential to accept.

\begin{observation}[Heterogeneity of kind]
\label{obs:kind}
Among the public resources a biological question routinely crosses, one finds
at least the following four kinds:
\begin{enumerate}[label=(\roman*),leftmargin=2.6em]
\item a SPARQL protocol endpoint accepting arbitrary graph patterns;
\item a resource exposing both a SPARQL endpoint and an independent REST
interface, with different coverage and different identifier conventions on the
two;
\item a REST interface returning flat delimited text, with a fixed finite set
of operations and no query language whatsoever;
\item an artefact --- an ontology file, a mapping table --- distributed for
local loading, whose query capability is whatever the loader provides.
\end{enumerate}
\end{observation}

Kind (iii) is the one that forecloses the obvious designs. A resource whose
entire query surface is a fixed family of request shapes such as
\texttt{/link/\{target\}/\{source\}} and \texttt{/get/\{id\}} does not admit a
compilation target for arbitrary patterns. It admits exactly the operations it
names. Any language that compiles to ``a query'' must therefore either exclude
such resources or define ``query'' so broadly that the word does no work. We
take the second horn and rename: what a source accepts is a \emph{request}, and
what shapes the request is a \emph{capability set}.

\begin{definition}[Capability vocabulary]
\label{def:capability}
Fix a finite vocabulary $\Feat$ of \emph{capabilities}. A capability names a
construct that a request may use. Throughout,
\[
\Feat = \{\textsf{pattern}, \textsf{path}, \textsf{bind}, \textsf{agg},
\textsf{filter}, \textsf{neg}, \textsf{order}, \textsf{lookup},
\textsf{link}, \textsf{batch}, \textsf{regex}\},
\]
where \textsf{pattern} is a conjunctive pattern over the source's vocabulary,
\textsf{path} a regular path expression, \textsf{bind} the inline supply of a
value set as input, \textsf{agg} an aggregate, \textsf{filter} a scalar
predicate evaluated at the source, \textsf{neg} negation, \textsf{order} a
sort, \textsf{lookup} retrieval of a record by primary identifier,
\textsf{link} retrieval of identifiers in a second namespace related to a given
identifier, \textsf{batch} the acceptance of more than one input identifier per
request, and \textsf{regex} a string match.
\end{definition}

The vocabulary is not canonical and nothing below depends on its particular
members; what the development requires is only that $\Feat$ be finite and that
capabilities occupy syntactically independent positions in each concrete
request language (\cref{asm:independence}).

\begin{definition}[Source]
\label{def:source}
A \emph{source} is a quadruple
$\Src = (\eta_\Src, \Capset(\Src), \Ext_\Src, \cost_\Src)$
where $\eta_\Src$ is an opaque endpoint handle, $\Capset(\Src) \subseteq \Feat$ is
the \emph{declared capability set}, $\Ext_\Src$ is the \emph{extraction
function} mapping a raw response to a result set, and
$\cost_\Src : \Nset \to \Rset_{>0}$ gives the cost of a request as a function
of its input cardinality.
\end{definition}

\begin{example}[Four sources]
\label{ex:sources}
A reaction knowledge base with a SPARQL endpoint declares
\[
\{\textsf{pattern},\textsf{path},\textsf{bind},\textsf{agg},\textsf{filter},
\textsf{neg},\textsf{order},\textsf{regex},\textsf{batch}\}.
\]
A flat-file pathway REST interface declares $\{\textsf{lookup},\textsf{link}\}$
and nothing else --- in particular not \textsf{filter}, so a plan wishing to
filter a pathway result must retrieve and filter locally, at a cost the plan
can see. A locally loaded ontology declares
$\{\textsf{pattern},\textsf{path},\textsf{filter}\}$ but not \textsf{batch},
since the loader admits one query at a time. A protein resource exposing both a
SPARQL endpoint and a REST interface is modelled as \emph{two} sources with
different capability sets, for the reason given in \cref{rem:not-union}.
\end{example}

\begin{remark}[A dual-interface resource is two sources]
\label{rem:not-union}
Modelling a resource with both SPARQL and REST access as a single source whose
capability set is the union of the two is unsound. The union asserts that some
request may use \textsf{path} and \textsf{link} together, and no request can:
the two capabilities belong to different interfaces, with different coverage
and often different identifier conventions. Capability sets are per-interface,
and a plan crossing the two interfaces of one resource crosses two sources and
pays two costs.
\end{remark}

\subsection{Results, and why the result model is the common one}

\begin{definition}[Result set]
\label{def:result}
A \emph{result set} over namespace $n$ is a finite set
\[
\Res \subseteq n \times (\mathcal{A} \rightharpoonup \mathcal{V})
\]
of pairs (identifier, partial attribute map), in which no identifier occurs
twice. We write $\idm(\Res) \subseteq n$ for the projection onto identifiers,
and $|\Res| = |\idm(\Res)|$.
\end{definition}

Making the identifier primary and the attributes secondary is the whole of the
portability argument, so it is worth establishing that the choice is available.

\begin{proposition}[The result model is common where the query model is not]
\label{prop:common-result}
Each of the four source kinds of \cref{obs:kind} admits an extraction function
into \cref{def:result} that is defined on every syntactically valid response.
\end{proposition}

\begin{proof}
For kind (i), and for the SPARQL interface of kind (ii), a response is a
solution sequence: a list of maps from projected variables to values. Designate
one projected variable as the identifier column and the remainder as
attributes; collapse duplicate identifiers by merging their attribute maps.
The construction is defined on every well-formed solution sequence with at
least one projected variable.

For the REST interfaces of kinds (ii) and (iii), a response is one of two
shapes. Under \textsf{lookup} it is a record keyed by the requested identifier,
yielding the singleton whose identifier is that key and whose attribute map is
the record's fields. Under \textsf{link} it is a delimited list of pairs,
yielding one element per line whose identifier is the second column and whose
attribute map records the first column as provenance. Both are defined on every
syntactically valid response.

For kind (iv), the loader's own result form is that of kind (i), and the
construction of the first paragraph applies. In each case the construction
consumes the entire response, so no information is discarded silently, though
information may be relocated into attributes.
\end{proof}

\begin{remark}[What \cref{prop:common-result} does not say]
\label{rem:extraction-lossy}
It does not say the extraction is lossless in the sense that matters. A
solution sequence with three projected variables and no functional dependency
on any one of them --- a genuine ternary relation --- is forced into
identifier-plus-attributes by choosing a key, and the choice loses the
multiplicity of the other two columns. A plan that needs an $n$-ary relation
not functionally determined by a key cannot express it in this model.
\Cref{sec:limits} records this as a real restriction rather than a formality,
and \cref{rem:nary-workaround} gives the standard, unsatisfying workaround.
\end{remark}

\begin{remark}[The workaround, and what it costs]
\label{rem:nary-workaround}
The standard response to the restriction of \cref{rem:extraction-lossy} is
reification: introduce a synthetic identifier for each tuple of the $n$-ary
relation and carry the original columns as attributes of it. The result is a
result set in the sense of \cref{def:result}, so the model accepts it, and
nothing is lost in the sense of information content.

What is lost is everything downstream. A synthetic identifier lies in no
namespace of any source, so it is outside the domain of every translation map
of \cref{def:map}; a reified step therefore cannot be the input of a
translation step, and the retention arithmetic of \cref{thm:retention} has
nothing to say about it. Nor can it be joined against a source on identity,
since no source will have minted the same synthetic key. In effect reification
buys expressiveness at the price of federation: the reified set is a terminal
value the plan may emit but may not carry across a boundary. We call the
workaround standard because it is what one does, and unsatisfying because a
language for federated retrieval that answers ``express it as something you
cannot federate'' has not answered.
\end{remark}

\begin{remark}[Disjoint namespaces are a modelling decision]
\label{rem:disjoint}
We assume namespaces are pairwise disjoint. In practice they are not: the same
lexical string may be a valid accession in two databases, and one database may
mint identifiers in another's namespace as aliases. The assumption is
discharged in the prototype by prefixing every identifier with its namespace
tag at extraction time, which makes disjointness true by construction at the
cost of making equality between an identifier and its own alias fail. That
failure is visible as a retention loss in \cref{sec:translation} rather than as
a silent join miss, which we regard as the better trade, but it is a trade.
\end{remark}

\subsection{Lowering, faithfulness, and honesty}

\begin{definition}[Abstract request]
\label{def:areq}
An \emph{abstract request} is a triple $\rho = (\phi, \Req(\rho), \beta)$ where
$\phi$ is a pattern over a declared vocabulary, $\Req(\rho) \subseteq \Feat$ is
the set of capabilities that $\phi$ uses, and $\beta$ is a possibly empty tuple
of result sets to be supplied as input. We write $\llbracket \rho \rrbracket_D$
for the \emph{denotation} of $\rho$ over a dataset $D$: the result set that
$\rho$ specifies, defined independently of any source.
\end{definition}

\begin{definition}[Lowering]
\label{def:lowering}
A \emph{lowering} for source $\Src$ is a partial function
$\Low_\Src : \rho \rightharpoonup \mathcal{W}_\Src$ into the source's concrete
request language. $\Low_\Src$ is \emph{faithful for} $f \in \Feat$ if for every
abstract request $\rho$ with $\Req(\rho) \subseteq \{f\} \cup \Capset(\Src)$ and
every dataset $D$ the source holds,
\[
\Ext_\Src\bigl(\eta_\Src(\Low_\Src(\rho))\bigr) \;=\; \llbracket \rho \rrbracket_D .
\]
\end{definition}

\begin{definition}[Honest declaration]
\label{def:honest}
$\Capset(\Src)$ is \emph{honest} if, for every $f \in \Feat$,
\[
f \in \Capset(\Src) \iff \Low_\Src \text{ is faithful for } f .
\]
\end{definition}

\begin{assumption}[Syntactic independence]
\label{asm:independence}
For each source $\Src$ there is an injection from $\Feat$ into a set of
pairwise disjoint syntactic positions of $\mathcal{W}_\Src$, such that the
concrete request realising a set of capabilities is the one filling exactly the
corresponding positions, and its denotation is the conjunction of the
denotations of the parts.
\end{assumption}

\Cref{asm:independence} holds for the concrete languages we target --- a
pattern, a path inside a pattern, a binding block, an aggregate wrapper, a
filter clause and a sort modifier are separate positions in SPARQL, and a REST
interface with a fixed operation family satisfies it degenerately because each
capability corresponds to a distinct operation. It fails for languages in which
two constructs interact semantically, and \cref{rem:independence-fails} records
one such interaction that is not hypothetical.

\begin{theorem}[Lowering is total and faithful exactly on contained requests]
\label{thm:lowering}
Let $\Src$ be a source whose declaration is honest and which satisfies
\cref{asm:independence}. Then for every abstract request $\rho$,
\[
\Low_\Src(\rho) \text{ is defined and faithful}
\quad\Longleftrightarrow\quad
\Req(\rho) \subseteq \Capset(\Src).
\]
\end{theorem}

\begin{proof}
($\Leftarrow$) Suppose $\Req(\rho) \subseteq \Capset(\Src)$. By
\cref{def:honest}, $\Low_\Src$ is faithful for each $f \in \Req(\rho)$
individually. By \cref{asm:independence} the capabilities in $\Req(\rho)$
occupy pairwise disjoint syntactic positions, so the concrete request filling
all of those positions is well formed and lies in $\mathcal{W}_\Src$; hence
$\Low_\Src(\rho)$ is defined. Again by \cref{asm:independence} its denotation
is the conjunction of the denotations of its parts, and by faithfulness for
each $f$ separately each part denotes what $\rho$'s corresponding component
specifies. The conjunction therefore equals $\llbracket \rho \rrbracket_D$, so
$\Low_\Src(\rho)$ is faithful.

($\Rightarrow$) Suppose $f \in \Req(\rho) \setminus \Capset(\Src)$ for some $f$.
By honesty, $\Low_\Src$ is not faithful for $f$: there exist an abstract
request $\rho'$ using $f$ and a dataset $D$ on which either $\Low_\Src(\rho')$
is undefined or the extracted response differs from
$\llbracket \rho' \rrbracket_D$. In the first case $\Low_\Src$ is undefined on
some request using $f$, and by \cref{asm:independence} the position that $f$
occupies is unfillable for $\Src$, hence $\Low_\Src(\rho)$ is undefined as
well. In the second case the position is fillable but its part denotes
something other than what $\rho$ specifies, and since by
\cref{asm:independence} the whole denotation is the conjunction over parts,
$\Low_\Src(\rho)$ is defined but unfaithful. Either way the left-hand side
fails.
\end{proof}

\begin{remark}[Honesty is an assumption, and the interesting one]
\label{rem:honesty-assumption}
\Cref{thm:lowering} is exactly as strong as \cref{def:honest}, and honesty is
asserted by the adapter author, not verified against the service. There are two
ways to violate it and they are not symmetric. \emph{Under}-declaration --- a
capability the source in fact supports faithfully, omitted from $\Capset(\Src)$ ---
is safe: it costs expressiveness, causes some plans to be refused that could
have been answered, and never produces a wrong answer. \emph{Over}-declaration
is unsound and invisible: the adapter claims \textsf{path}, the service accepts
path expressions but evaluates them under a different semantics or truncates
them, the lowering is defined, a response comes back, and it is not the
denotation. This is not hypothetical --- \cref{sec:interp} exhibits a case in
which a public service accepted a construct and returned a count differing from
the specified denotation by a factor of $198.5$. We therefore treat the
capability set as a claim to be minimised rather than a feature list to be
maximised, and record in \cref{sec:limits} that no mechanical check enforces
this, and that we know of no way to build one short of differential testing
against a reference implementation.
\end{remark}

\begin{proposition}[Under-declaration can be forced by the service, not the language]
\label{prop:under}
There exist a source $\Src$, a capability $f$, and a concrete language
$\mathcal{W}_\Src$ such that $\mathcal{W}_\Src$ supports $f$ in full
generality, $\Low_\Src$ emits well-formed requests using $f$, every such
request receives a well-formed response, and honesty nonetheless requires
$f \notin \Capset(\Src)$.
\end{proposition}

\begin{proof}
Take $f = \textsf{agg}$ and let $\Src$ be a SPARQL endpoint operating under a
server-side cap $M$ on the number of solutions it will materialise. Aggregation
is expressible in SPARQL in full generality, so $\mathcal{W}_\Src$ supports
$f$. Let $\Low_\Src$ emit a counting aggregate over a pattern $\phi$ whose
unaggregated solution sequence has cardinality $N > M$ on the dataset $D$ the
source holds. The endpoint materialises $M$ solutions, computes the aggregate
over those, and returns a well-formed response carrying the number $M$. But
$\llbracket \rho \rrbracket_D = N \neq M$. So $\Low_\Src$ is defined on $\rho$
and unfaithful for $f$, and by \cref{def:honest} honesty forces
$f \notin \Capset(\Src)$ even though the target language supports $f$ and the
request succeeded.
\end{proof}

\Cref{prop:under} is the reason \cref{def:honest} is phrased in terms of the
lowering and the service jointly rather than in terms of the target language's
grammar. Faithfulness is not a property of a language; it is a property of a
construct, a lowering, and a deployment. A capability table derived from a
specification is therefore not a capability declaration, and treating it as one
is precisely the over-declaration error of \cref{rem:honesty-assumption}.

\begin{remark}[Where \cref{asm:independence} fails]
\label{rem:independence-fails}
The assumption fails when two constructs interact. In SPARQL, a path expression
inside a pattern and an inline value block are formally independent positions,
but the \emph{expansion} of an abbreviated object list happens during
translation to the abstract algebra, before evaluation, and interacts with
path expressions in a way that surprises authors --- as \cref{sec:interp}
documents. Where independence fails, \cref{thm:lowering} does not apply and the
lowering must be verified for the interacting pair jointly. The prototype
handles this by never emitting the abbreviated form (\cref{cons:longhand}),
which restores independence by removing one of the interacting constructs from
the emitted language altogether.
\end{remark}

\section{The intermediate representation}
\label{sec:ir}

\subsection{Steps}

\begin{definition}[Step]
\label{def:step}
A \emph{step} is a tuple
$\Step = (x, \Src, \rho, \beta, b)$
where $x$ is the variable the step binds, $\Src$ is a source, $\rho$ is an
abstract request, $\beta = (y_1,\dots,y_k)$ is a tuple of variables whose
values are supplied to $\rho$ as input, and $b \in \Rset_{>0} \cup \{\infty\}$
is the step's budget annotation.
\end{definition}

The shape $\Step = \text{source} \times \text{request} \times \text{binding}$
is the whole of the intermediate representation, and it is the reason a source
adapter is small: of the five components, only $\Low_\Src$ and $\Ext_\Src$ are
source-specific, and both are functions of the abstract request alone.

\begin{definition}[Plan]
\label{def:plan}
A \emph{plan} $\Plan$ is a finite sequence $\Step_1,\dots,\Step_m$ of steps
together with a total budget $\Bud$, such that the variable bound by $\Step_i$
is distinct from that bound by $\Step_j$ for $i \neq j$, and every variable in
$\beta_i$ is bound by some $\Step_j$ with $j < i$. The \emph{dependency graph}
$G(\Plan)$ has the steps as vertices and an edge $\Step_j \to \Step_i$ whenever
$y \in \beta_i$ is bound by $\Step_j$.
\end{definition}

By construction $G(\Plan)$ is a directed acyclic graph and the sequence order
is one of its topological orders. Nothing below depends on which one.

\begin{example}[The motivating question as a plan]
\label{ex:plan}
\begin{lstlisting}[caption={The question of \cref{sec:intro} as a plan. Compare \cref{lst:script}: no query text, no batching constant, explicit budget, explicit failure handling.},label=lst:plan]
plan enzymes_in_shared_pathways {
  budget 400 requests

  let acids = from CHEBI
      ask descendants_of("CHEBI:35238")
      within 20

  let rxns  = from RHEA
      ask reactions_consuming(?c)
      with ?c in acids
      within 120
      else fail unresolved

  let ecs   = from RHEA
      ask enzyme_of(?r)
      with ?r in rxns
      within 60

  let kids  = map ecs via ec_to_kegg
      expect partial 0.6

  let paths = from KEGG
      ask link("pathway", ?k)
      with ?k in kids
      within 200
      when starved emit partial

  emit paths with provenance
}
\end{lstlisting}
The \texttt{expect partial 0.6} annotation on the translation step is not
decoration: it is the assertion checked by \cref{def:retention-check}, and a
plan whose translation retains less than the asserted fraction stops rather
than proceeding on a silently decimated input --- which is defect (D3) of
\cref{sec:intro}, caught at the step that causes it rather than at the step that
suffers it.
\end{example}

\subsection{The static capability check}

\begin{definition}[Well-capability]
\label{def:wellcap}
A step $\Step = (x,\Src,\rho,\beta,b)$ is \emph{well-capability} if
$\Req(\rho) \subseteq \Capset(\Src)$, and additionally
$\textsf{batch} \in \Capset(\Src)$ whenever $\beta \neq ()$ and the step is
annotated for batched execution. A plan is well-capability if every step is.
\end{definition}

\begin{theorem}[Static decidability of compilability]
\label{thm:static}
Let $\Plan$ be a plan over sources with honest declarations satisfying
\cref{asm:independence}. Then:
\begin{enumerate}[label=(\alph*),leftmargin=2.4em]
\item whether $\Plan$ is well-capability is decidable in time
$O\bigl(m \cdot |\Feat|\bigr)$, where $m$ is the number of steps, without
issuing any request;
\item if $\Plan$ is well-capability then every $\Low_{\Src_i}(\rho_i)$ is
defined and faithful;
\item if $\Plan$ is not well-capability then there is a step $\Step_i$ and a
capability $f \in \Req(\rho_i) \setminus \Capset(\Src_i)$ such that no execution
of $\Plan$ can produce a faithful result for $\Step_i$, and the pair
$(\Step_i, f)$ is computable within the same bound.
\end{enumerate}
\end{theorem}

\begin{proof}
(a) $\Req(\rho_i)$ is determined syntactically from $\rho_i$ by structural
recursion over the abstract request, and $\Capset(\Src_i)$ is a declared finite
set; the containment test is a subset test over a fixed finite vocabulary,
costing $O(|\Feat|)$ per step with bitset representation. There are $m$ steps
and no step's test depends on any other step's result, so the total is
$O(m|\Feat|)$. Nothing in the test consults $\eta_{\Src_i}$, so no request is
issued.

(b) Immediate from \cref{thm:lowering} applied stepwise, the hypotheses of
which hold by assumption.

(c) If $\Plan$ is not well-capability then some step fails the containment
test; take the least such $i$ in the sequence order and any witness
$f \in \Req(\rho_i)\setminus\Capset(\Src_i)$, both produced by the same scan. By
\cref{thm:lowering}, $\Low_{\Src_i}(\rho_i)$ is either undefined or unfaithful.
In the first case no request is ever issued for $\Step_i$ and it has no result.
In the second, a result is produced but differs from
$\llbracket \rho_i \rrbracket_D$ on some dataset, so it is not faithful. Since
the executor's behaviour at $\Step_i$ depends on the source only through
$\Low_{\Src_i}$ and $\Ext_{\Src_i}$, no scheduling choice, retry, or budget
allocation changes this.
\end{proof}

\begin{corollary}[Refusal before contact]
\label{cor:refuse-before-contact}
A plan that is not well-capability can be refused, with the offending
(step, capability) pair named, before any network request is issued.
\end{corollary}

\begin{proof}
Immediate from \cref{thm:static}(a) and (c): the test is decidable without
issuing a request, and on failure it computes the witness.
\end{proof}

\Cref{cor:refuse-before-contact} is \cref{prin:refusal} made operational, and
it is the concrete payoff of separating $\Capset$ from $\Ext$ in
\cref{def:source}. It is worth being precise about what it buys. It does not
guarantee that a well-capability plan succeeds --- \cref{sec:verdicts} lists
four ways it can fail with every capability present, and
\cref{prop:necessary-not-sufficient} adds a fifth. It guarantees only that one
particular family of failures, the family in which a plan asks a source for
something it cannot do, is converted from a runtime symptom into a
compile-time diagnosis naming its own cause.

\begin{proposition}[The check is not conservative in the useless direction]
\label{prop:check-tight}
There is no plan that is refused by \cref{def:wellcap} and that could
nonetheless be executed faithfully against sources with honest declarations.
\end{proposition}

\begin{proof}
Suppose $\Plan$ is refused. Then some step has
$f \in \Req(\rho_i)\setminus\Capset(\Src_i)$, and by \cref{thm:static}(c) no
execution produces a faithful result for that step. So the refusal excludes no
faithfully executable plan.
\end{proof}

\Cref{prop:check-tight} is the converse guarantee to
\cref{cor:refuse-before-contact} and is what distinguishes a capability check
from a conservative approximation: the check refuses exactly the plans that
cannot work, given honest declarations. Under \emph{dishonest} declarations it
inherits the dishonesty in both directions, which is
\cref{rem:honesty-assumption} again and is the reason that remark, rather than
this proposition, is the honest summary of what has been established.

%==============================================================================
\part{Verdicts}
%==============================================================================

\section{The verdict algebra}
\label{sec:verdicts}

\subsection{Three independent predicates}

Before naming the verdicts, we isolate what a verdict is about. Fix a step
$\Step = (x, \Src, \rho, \beta, b)$ executed against a source holding dataset
$D$, with the values of $\beta$ already computed. Three predicates apply, and
the burden of this subsection is that they are logically independent.

\begin{definition}[The three predicates]
\label{def:three}
\leavevmode
\begin{itemize}[leftmargin=2.4em]
\item $\mathrm{Exp}(\Step)$ holds iff the abstract request $\rho$ is
expressible against $\Src$: formally, $\Req(\rho) \subseteq \Capset(\Src)$.
\item $\mathrm{Der}(\Step)$ holds iff
$\llbracket \rho \rrbracket_D \neq \emptyset$ on the dataset the source in fact
holds, with the supplied bindings $\beta$.
\item $\mathrm{Ans}_b(\Step)$ holds iff the concrete request completes within
the step's budget $b$.
\end{itemize}
\end{definition}

\begin{proposition}[Independence]
\label{prop:independence}
The three predicates of \cref{def:three} are logically independent: all eight
truth assignments are realisable by some (step, source, dataset, budget).
\end{proposition}

\begin{proof}
We exhibit the constructions; each varies one predicate holding the others
fixed, and the eight cases follow by composing them.

To vary $\mathrm{Exp}$ alone: take a request using \textsf{regex} and swap the
source between one declaring \textsf{regex} and one not. Neither the dataset
nor the budget changes, so $\mathrm{Der}$ and $\mathrm{Ans}_b$ are untouched on
the branch where the request is issued at all; on the branch where it is not,
$\mathrm{Der}$ is still a property of $D$ and $\rho$, both fixed, and
$\mathrm{Ans}_b$ is vacuously assigned by convention.

To vary $\mathrm{Der}$ alone: fix the source and request and vary the dataset
between one containing a witness for $\phi$ and one not. Capability
declarations do not depend on $D$, so $\mathrm{Exp}$ is unchanged; a request
that completes on one finite dataset completes on the other for a budget
exceeding both costs, so $\mathrm{Ans}_b$ can be held.

To vary $\mathrm{Ans}_b$ alone: fix source, request and dataset, and vary $b$
across the cost of the request. Neither $\Capset(\Src)$ nor
$\llbracket\rho\rrbracket_D$ depends on $b$, so the other two are unchanged.

Since each predicate is separately variable with the other two held fixed, and
each is two-valued, all eight assignments are realisable.
\end{proof}

\Cref{prop:independence} is the reason a single boolean cannot carry the
outcome of a step: it must report a point in an eight-element space, and one
bit indexes two. The verdict set below is the coarsening of that space which
distinguishes the cases calling for different repairs.

\begin{remark}[Only one predicate carries a clock]
\label{rem:clock}
$\mathrm{Exp}$ is a property of the request and the declaration;
$\mathrm{Der}$ is a property of the request and the dataset; only
$\mathrm{Ans}_b$ mentions the budget. A system reporting a single outcome
cannot say which of the three it means, and a practitioner who reruns a failing
step with more time is implicitly betting that the failure was in the third.
The bet is unfounded two-thirds of the time, and \cref{lst:script} provides no
way to check it.
\end{remark}

\subsection{Six verdicts}

\begin{definition}[Verdict set]
\label{def:verdicts}
$V = \{\vans, \vempty, \vsurf, \vtime, \vref, \vstarve\}$, assigned to a step
by the following rules, applied in order:
\begin{enumerate}[label=(R\arabic*),leftmargin=3em]
\item If some $y \in \beta$ has a verdict other than $\vans$, or has verdict
$\vans$ with a payload whose retention fell below the step's declared
expectation, then $\vstarve$; the diagnosis names the offending predecessor.
\item Else if $\Req(\rho) \not\subseteq \Capset(\Src)$ then $\vsurf$; the
diagnosis names $\Req(\rho)\setminus\Capset(\Src)$.
\item Else if the plan's remaining budget is below $\cost_\Src$ evaluated at
the input cardinality, then $\vref$; the diagnosis names the shortfall.
\item Else issue the request. If it does not complete within $b$ then $\vtime$;
the diagnosis names $b$ and the elapsed cost.
\item Else if the extracted result set is empty then $\vempty$.
\item Else $\vans$, with the result set as payload.
\end{enumerate}
\end{definition}

The ordering of (R1)--(R6) is not arbitrary and is not merely an implementation
convenience: it is what makes the verdicts \emph{attributive}. Each rule fires
only when every earlier cause has been excluded, so a step reporting $\vempty$
has established that its inputs were adequate, its capabilities sufficient, its
budget available and its request completed --- and therefore that the emptiness
is a fact about the data rather than about the machinery. No rule below (R6)
can say that.

\begin{definition}[Blocker]
\label{def:blocker}
The \emph{blocker} of a non-$\vans$ verdict is
\[
\mathrm{blk}(\vsurf) = \bmdl,\quad
\mathrm{blk}(\vtime) = \beng,\quad
\mathrm{blk}(\vref) = \bbud,\quad
\mathrm{blk}(\vstarve) = \bcorp,
\]
and $\vempty$ has no blocker.
\end{definition}

\begin{remark}[$\vempty$ has no blocker on purpose]
Assigning a blocker to $\vempty$ would assert that something obstructed the
step, and (R6) fires exactly when nothing did. A verdict of $\vempty$ is a
successful measurement of an empty set. Conflating it with an obstruction is
defect (D2) of \cref{sec:intro} in its purest form, and the discipline that
prevents it is that the blocker map is partial.
\end{remark}

\subsection{Distinguishability, and the one-bit interface}

\begin{theorem}[Six-way distinguishability]
\label{thm:six}
For each pair $v \neq v'$ in $V$ there exist a plan, two source
configurations agreeing on everything the plan can observe except as noted, and
a step whose verdict is $v$ in the first and $v'$ in the second, such that the
executor of \cref{def:verdicts} reports the difference. Consequently the six
verdicts are pairwise distinguishable by a plan executor.
\end{theorem}

\begin{proof}
It suffices to realise each verdict by a configuration differing from a common
baseline in exactly one respect, since then any two verdicts are separated by
the pair of configurations realising them.

Take as baseline a step whose predecessors all returned $\vans$ at full
retention, whose source declares every capability the request uses, whose
budget exceeds the request cost, and whose dataset contains a witness. By
(R1)--(R6) the baseline verdict is $\vans$.

$\vempty$: remove the witness from $D$, changing nothing else. (R1)--(R4) still
pass; (R5) fires.

$\vtime$: restore the witness and set $b$ below the request cost. (R1)--(R3)
pass, (R4) fires.

$\vref$: restore $b$ and set the plan's remaining budget below
$\cost_\Src$. (R1)--(R2) pass, (R3) fires.

$\vsurf$: restore the budget and remove one capability of $\Req(\rho)$ from
$\Capset(\Src)$. (R1) passes, (R2) fires.

$\vstarve$: restore the capability and set a predecessor's verdict to any
non-$\vans$ value, or its retention below the declared expectation. (R1) fires.

Each configuration differs from the baseline in exactly one respect and yields
a distinct verdict, and the executor computes the verdict by the stated rules,
so it reports each. The six are therefore pairwise distinguishable.
\end{proof}

\begin{corollary}[A row-returning interface conflates at least four verdicts]
\label{cor:onebit}
Let $I$ be an interface whose response to a request is a result set, with
failure signalled by returning the empty set. Then $I$ cannot distinguish
$\vempty$, $\vtime$, $\vsurf$ and $\vstarve$ in any of the four configurations
constructed in the proof of \cref{thm:six}. No discipline imposed on callers of
$I$ repairs this.
\end{corollary}

\begin{proof}
In each of the four configurations the value $I$ returns is the empty result
set: in the $\vempty$ configuration because the denotation is empty; in the
$\vtime$ configuration because the request did not complete and $I$ has no
other channel; in the $\vsurf$ configuration because the request was not issued
or was issued and degraded; in the $\vstarve$ configuration because the inputs
were empty and a request over empty inputs has empty denotation. The four
responses are equal as values. A caller's behaviour is a function of the values
it receives, so no caller distinguishes them, and no discipline on callers
changes the values. Hence the conflation is a property of $I$, not of its use.
\end{proof}

\Cref{cor:onebit} is the formal content of defect (D2), and it is worth being
explicit about its scope. It does not say that existing systems are badly
engineered. It says that the interface contract they expose --- rows, with
emptiness overloaded as failure --- has insufficient capacity, and that a
system built on it cannot recover the distinction downstream however carefully
it is written. The repair is not better error handling; it is a wider return
type.

\subsection{Starvation is the characteristic federated failure}

The verdict $\vstarve$ deserves separate treatment, because it is the one with
no single-database analogue, and because it is not a corner case.

\begin{proposition}[Starvation is unreachable in the closed setting and reachable in the federated one]
\label{prop:starve-reachable}
Consider the restriction of \cref{def:verdicts} to plans of a single step
against a fixed corpus --- the closed setting. Then:
\begin{enumerate}[label=(\alph*),leftmargin=2.4em]
\item no execution assigns $\vstarve$, and consequently the blocker $\bcorp$
is emitted by no rule;
\item in the federated setting with $m \geq 2$ steps, $\vstarve$ is assigned on
a set of configurations of positive measure in the following sense: for every
step $\Step_i$ with $\beta_i \neq ()$ and every predecessor $\Step_j$ feeding
it, there is a dataset perturbation at $\Src_j$ alone that causes $\Step_i$ to
report $\vstarve$ while $\Step_i$'s own source, capabilities and budget are
untouched.
\end{enumerate}
\end{proposition}

\begin{proof}
(a) In a single-step plan, $\beta = ()$, so the antecedent of (R1) is vacuously
false and (R1) never fires. Since (R1) is the unique rule assigning $\vstarve$,
the verdict is unreachable; and since $\bcorp$ is by \cref{def:blocker} the
image of $\vstarve$ alone, no blocker $\bcorp$ is emitted. The corpus is fixed
and given, so its inadequacy manifests as $\vempty$ under (R5) --- an empty
answer, correctly reported as such, with no obstruction named.

(b) Let $\Step_j \to \Step_i$ be an edge of $G(\Plan)$ with $y \in \beta_i$
bound by $\Step_j$. Perturb the dataset at $\Src_j$ by deleting the witnesses
for $\rho_j$. Then $\Step_j$ reports $\vempty$ by (R5); its capability set,
budget and lowering are unchanged, so no other verdict applies to it. At
$\Step_i$, the antecedent of (R1) now holds, since $y$ has a verdict other than
$\vans$; (R1) fires and $\Step_i$ reports $\vstarve$ naming $\Step_j$. Nothing
at $\Src_i$ was altered.
\end{proof}

\begin{remark}[Why this is the characteristic case]
\label{rem:starve-characteristic}
In the closed setting the corpus is an input: if it lacks the assertions a
question needs, that is a fact about the world one reports and does not
diagnose. In federation the corpus of step $i$ is not an input but an
\emph{artefact of steps $1..i-1$}, and its inadequacy is therefore a
diagnosable event with a locatable cause. \Cref{prop:starve-reachable} says
that federation makes a previously unreachable diagnosis reachable, and
\cref{prop:starve-common} below says it is the modal outcome once translation
enters. Read together with \cref{cor:onebit}: the failure mode that federation
newly creates is exactly one of the four that the conventional interface cannot
see.
\end{remark}

\begin{proposition}[Starvation dominates once translation enters]
\label{prop:starve-common}
Let $\Plan$ contain a translation step whose retention is $r < 1$
(\cref{def:retention}), feeding a step whose per-input hit probability is $p$
independently across inputs. Then for input cardinality $N$ the probability
that the downstream step returns a non-empty payload is
$1 - (1-p)^{rN}$, whereas absent the translation it is $1-(1-p)^{N}$. For $p$
small and $rN$ small the ratio of failure probabilities approaches
$e^{-prN}/e^{-pN} = e^{p N (1-r)}$, so the excess failure attributable to
translation grows exponentially in the lost fraction $1-r$ times the input
size.
\end{proposition}

\begin{proof}
The translation delivers $rN$ inputs by \cref{def:retention}. Under
independence the downstream step misses on all of them with probability
$(1-p)^{rN}$, giving the stated success probability; the untranslated case is
$r=1$. Taking the ratio of the two failure probabilities and using
$(1-p)^{k} = e^{k\ln(1-p)} \approx e^{-pk}$ for small $p$ yields
$e^{-prN + pN} = e^{pN(1-r)}$.
\end{proof}

The independence hypothesis in \cref{prop:starve-common} is false in practice
--- identifiers that fail to translate are correlated with identifiers that are
poorly annotated downstream, so the true excess is larger than the bound and
not smaller. We state the independent case because it is the one we can prove,
and record the direction of the bias rather than claiming the bound is tight.

\subsection{Verdict propagation}

\begin{definition}[Propagation]
\label{def:propagation}
A plan's verdict is $\vans$ if every emitted step is $\vans$; otherwise it is
the verdict of the earliest non-$\vans$ step in the sequence order, together
with the transitive closure of blame along $G(\Plan)$.
\end{definition}

\begin{proposition}[Blame is well defined and terminates at a non-starved step]
\label{prop:blame}
For any plan and execution, following the diagnosis of a $\vstarve$ verdict
from step to named predecessor terminates, and terminates at a step whose
verdict is not $\vstarve$.
\end{proposition}

\begin{proof}
By \cref{def:plan}, every variable in $\beta_i$ is bound by some $\Step_j$ with
$j < i$, so the named predecessor is strictly earlier in the sequence order.
The chain of blame is therefore strictly decreasing in a well-founded order on
a finite set and must terminate. It terminates at a step which is not
$\vstarve$, because a $\vstarve$ step always names a predecessor by (R1) and so
is never terminal; and the first step of the plan has $\beta_1 = ()$, so (R1)
cannot fire on it.
\end{proof}

\Cref{prop:blame} is what makes a $\vstarve$ report actionable rather than
circular: the analyst follows the chain and arrives at a step whose failure is
attributed to a capability, a budget, a timeout, or a genuinely empty dataset
--- one of the four things one can actually do something about. This is defect
(D3) resolved, and it is resolved by the structure of the verdict algebra
rather than by a convention that callers are asked to observe.

%==============================================================================
\part{Translation}
%==============================================================================

\section{The arithmetic of partial identifier translation}
\label{sec:translation}

Every federated plan that crosses a database boundary crosses a namespace
boundary, and the crossing is performed by a cross-reference map. These maps
are the least examined component of federated retrieval and the one where the
losses actually accumulate. This section treats them as first-class objects of
the language and derives what a plan can and cannot know about a chain of them.

\subsection{Maps are partial, non-injective, and non-confluent}

\begin{definition}[Translation map]
\label{def:map}
A \emph{translation map} from namespace $n$ to namespace $n'$ is a relation
$\mu \subseteq n \times n'$. We write $\mu(S) = \{ v : u \in S,\ (u,v) \in \mu\}$
for the image of a set, $\operatorname{dom}\mu = \{u : \exists v,\ (u,v) \in \mu\}$,
and call $\mu$ \emph{total} if $\operatorname{dom}\mu = n$, \emph{functional}
if each $u$ has at most one image, and \emph{injective} if each $v$ has at most
one preimage.
\end{definition}

Real cross-reference maps are none of the three. A chemical entity may have no
counterpart in a pathway database (partiality); a single ontology class may
correspond to several database compounds differing only in protonation state
(non-functionality); and several distinct ontology classes may map to one
database compound that does not resolve the distinction (non-injectivity).

\begin{definition}[Retention]
\label{def:retention}
The \emph{retention} of $\mu$ on a set $S$ is
\[
r_\mu(S) \;=\; \frac{|S \cap \operatorname{dom}\mu|}{|S|}
\qquad (S \neq \emptyset),
\]
the fraction of $S$ that survives translation, and $r_\mu(\emptyset) = 1$ by
convention. The \emph{amplification} is
$a_\mu(S) = |\mu(S)| / |S \cap \operatorname{dom}\mu|$ when the denominator is
nonzero.
\end{definition}

Retention and amplification are independent: a map may retain everything and
amplify (each input has several images), retain little and not amplify, or
both. Reporting only the output cardinality $|\mu(S)|$ conflates them, which is
why \cref{def:retention} separates them and why the prototype records both.

\begin{proposition}[Cardinality alone is uninformative]
\label{prop:cardinality-uninformative}
For any $N \geq 1$ and any output size $M \geq 1$ there exist maps $\mu_1,\mu_2$
and a set $S$ with $|S| = N$ such that
$|\mu_1(S)| = |\mu_2(S)| = M$ while $r_{\mu_1}(S) = 1$ and
$r_{\mu_2}(S) = 1/N$.
\end{proposition}

\begin{proof}
Let $S = \{u_1,\dots,u_N\}$. For $\mu_1$, distribute $M$ images across all $N$
inputs so that every input has at least one when $M \geq N$, and when $M < N$
let the inputs share images so that all $N$ lie in the domain and the image has
$M$ elements; then $r_{\mu_1}(S)=1$. For $\mu_2$, let
$\operatorname{dom}\mu_2 \cap S = \{u_1\}$ and let $\mu_2(u_1)$ have $M$
elements; then $|\mu_2(S)| = M$ and $r_{\mu_2}(S) = 1/N$. The output
cardinalities agree and the retentions differ by a factor of $N$.
\end{proof}

\Cref{prop:cardinality-uninformative} is the reason a plan cannot detect a
translation failure by looking at how much came out. In \cref{lst:script} the
translation is line~7 and its output size is the only thing the script observes
about it; by the proposition, that observation carries no information about
whether the translation lost anything. This is defect (D3), and it is not
fixable by checking that the output is non-empty.

\subsection{Composition}

\begin{theorem}[Retention of a chain]
\label{thm:retention}
Let $\mu_1 : n_0 \to n_1, \dots, \mu_k : n_{k-1} \to n_k$ and let $S_0$ be
finite and non-empty, with $S_i = \mu_i(S_{i-1})$. Then
\begin{enumerate}[label=(\alph*),leftmargin=2.4em]
\item the end-to-end retention factorises multiplicatively along the realised
sets,
\[
\frac{|S_k|}{|S_0|}
\;=\;
\prod_{i=1}^{k} r_{\mu_i}(S_{i-1})\, a_{\mu_i}(S_{i-1}),
\]
with the convention that the product is $0$ if any $S_{i-1}$ is empty;
\item if each $\mu_i$ is injective on the realised sets, the \emph{surviving
fraction} --- the fraction of $S_0$ having at least one image in $S_k$ ---
satisfies
\[
\rho_{1..k}(S_0) \;\leq\; \min_{1 \leq i \leq k} r_{\mu_i}(S_{i-1}),
\]
and the bound is attained when each $\mu_i$ is in addition functional. Without
injectivity the inequality fails, and $\rho_{1..k}(S_0) \leq r_{\mu_1}(S_0)$ is
all that survives (\cref{rem:injectivity-needed});
\item under the same injectivity hypothesis,
$\rho_{1..k}(S_0) \geq 1 - \sum_{i=1}^{k}\bigl(1 - r_{\mu_i}(S_{i-1})\bigr)$,
and the bound is attained when the loss sets are pairwise disjoint under the
realised correspondence. Injectivity is again load-bearing: without it a single
lost element of $S_{i-1}$ can extinguish several survivors of $S_0$ at once, and
the bound fails (\cref{rem:injectivity-needed}).
\end{enumerate}
\end{theorem}

\begin{proof}
(a) By \cref{def:retention}, $|S_i| = |\mu_i(S_{i-1})| =
a_{\mu_i}(S_{i-1}) \cdot |S_{i-1} \cap \operatorname{dom}\mu_i| =
a_{\mu_i}(S_{i-1}) r_{\mu_i}(S_{i-1}) |S_{i-1}|$ whenever $S_{i-1}$ is
non-empty. Telescoping over $i = 1..k$ and dividing by $|S_0|$ gives the
product. If some $S_{i-1}$ is empty then $S_k$ is empty and both sides are
zero.

(b) An element of $S_0$ survives to $S_k$ only if it lies in
$\operatorname{dom}\mu_1$, and its image lies in $\operatorname{dom}\mu_2$, and
so on. Let $A \subseteq S_0$ be the surviving set and fix a stage $i$. Each
$u \in A$ has at least one trajectory reaching a surviving element of
$S_{i-1}$, so the map sending $u$ to one such element is a function
$A \to S_{i-1} \cap \operatorname{dom}\mu_i$. If every $\mu_j$ is injective on
the realised sets this function is itself injective, whence
$|A| \leq r_{\mu_i}(S_{i-1})|S_{i-1}| \leq r_{\mu_i}(S_{i-1})|S_0|$ and the
minimum bound follows. Adding functionality gives each element exactly one
trajectory, and taking all losses to occur at the single stage achieving the
minimum attains the bound. The case $i=1$ needs no hypothesis, since survival
requires membership of $\operatorname{dom}\mu_1$ and distinct elements of $S_0$
are distinct there; this is the residual bound
$\rho_{1..k}(S_0) \leq r_{\mu_1}(S_0)$.

(c) The set of elements of $S_0$ that fail to survive is contained in the union
over $i$ of the sets $D_i \subseteq S_0$ whose trajectory dies at stage $i$.
A trajectory dies at stage $i$ by reaching an element of $S_{i-1}$ outside
$\operatorname{dom}\mu_i$, and there are $(1-r_{\mu_i}(S_{i-1}))|S_{i-1}|$ such
elements. Under injectivity on the realised sets, distinct members of $D_i$
reach distinct elements of $S_{i-1}$ --- the same injection constructed in (b)
--- so $|D_i| \leq (1-r_{\mu_i}(S_{i-1}))|S_{i-1}| \leq
(1-r_{\mu_i}(S_{i-1}))|S_0|$, the last step because injectivity also forces
$|S_{i-1}| \leq |S_0|$. A union bound over the $k$ stages gives the stated
lower bound, with equality when the $D_i$ are pairwise disjoint. Without
injectivity $|D_i|$ is not controlled by $|S_{i-1}|$ at all, since one dead
element of $S_{i-1}$ may be the sole continuation of arbitrarily many members
of $S_0$.
\end{proof}

\begin{remark}[Injectivity is a hypothesis of (b) and (c), not a nicety]
\label{rem:injectivity-needed}
The hypothesis in \cref{thm:retention}(b) and (c) cannot be dropped. Take
$S_0 = \{u_1,\dots,u_9\}$ with $\mu_1$ retaining seven of them, sending $u_9$
to $\{t_9\}$ and $u_{10}$ to $\{t_9, t_{10}\}$ --- non-injective, but with
amplification $a_{\mu_1}(S_0) = 1$, so the collision is invisible in the
factorisation of (a). Let $\mu_2$ retain three of the seven, among them $t_9$.
Then $r_{\mu_1} = 7/9$, $r_{\mu_2} = 3/7$, and four elements of $S_0$ survive,
giving $\rho = 4/9 > 3/7 = \min_i r_{\mu_i}$. Two survivors in $S_0$ are riding
one survivor in $S_1$, so the surviving set upstream is larger than the
surviving set downstream and the minimum is not an upper bound for it.

The failure is non-injectivity specifically, not amplification: $a_{\mu_1} = 1$
throughout the example. A cross-reference table that maps two source entries
onto one target entry --- routine wherever one resource merges what another
distinguishes --- is exactly this configuration, so the hypothesis is a real
restriction on which chains the two bounds may be applied to, and a plan
reporting coverage should measure $\rho$ rather than bound it.

The lower bound (c) fails for the same reason, by a different route. Its union
bound counts elements of $S_0$ whose trajectory dies at stage $i$ by counting
the elements of $S_{i-1}$ that lie outside $\operatorname{dom}\mu_i$; when
several survivors of $S_0$ have been merged onto one element of $S_{i-1}$, a
single such loss extinguishes all of them, and the count understates the
damage. A two-stage instance: $|S_0| = 8$ with realised sizes
$|S_1| = 6$, $|S_2| = 5$ and stage retentions $3/4$ and $5/6$, so (c) predicts
$\rho \geq 1 - 1/4 - 1/6 = 7/12$, while the measured surviving fraction is
$1/2$. The one element lost at stage two carried two survivors, costing $2/8$
where the bound had budgeted $1/6$. Note that the na\"ive repair ---
requiring $|S_{i-1}| \leq |S_0|$ --- does not help, since it already holds here.

Sampling two-stage chains over eight elements with each element retained
independently with probability $0.8$ and given up to three images, and keeping
only the non-injective draws, (b) fails for $33.9\%$ of chains and (c) for
$2.0\%$ ($n = 5000$). Under injective draws neither fails, as the theorem now
requires.
\end{remark}

\begin{remark}[The two bounds are far apart, and the gap is the point]
\label{rem:bounds-gap}
For $k=3$ and $r_{\mu_i} = 0.8$ throughout, on an injective chain, (b) gives
$\rho \leq 0.8$ and (c) gives $\rho \geq 0.4$. The true value depends on
whether the three maps lose \emph{the same} entities or \emph{different} ones,
and no amount of pairwise coverage data determines which. A plan that needs to
know its own coverage must therefore measure the chain, not compute it from
published per-map figures --- which is what \cref{def:retention-check} makes it
do, and what \cref{prop:pairwise-insufficient} says it cannot avoid.
\end{remark}

\begin{proposition}[Pairwise coverage does not determine chain coverage]
\label{prop:pairwise-insufficient}
There exist two families $(\mu_1,\mu_2)$ and $(\mu_1',\mu_2')$ with
$r_{\mu_i} = r_{\mu_i'}$ for $i=1,2$ on a common $S_0$, whose end-to-end
surviving fractions differ.
\end{proposition}

\begin{proof}
Let $S_0 = \{a,b,c,d\}$, and let both first maps be the identity on
$\{a,b,c\}$ and undefined on $d$, so $r_{\mu_1} = r_{\mu_1'} = 3/4$. Let
$\mu_2$ be the identity on $\{a,b\}$ and undefined on $c$, and let $\mu_2'$ be
the identity on $\{b,c\}$ and undefined on $a$; both have retention $2/3$ on
$S_1 = \{a,b,c\}$. The chain through $\mu_2$ retains $\{a,b\}$, a surviving
fraction of $2/4$; the chain through $\mu_2'$ retains $\{b,c\}$, also $2/4$. To
separate them, extend $\mu_2'$ to be defined on $a$ as well and undefined on
both $b$ and $c$: then $r_{\mu_2'} = 1/3 \neq 2/3$, so instead take
$S_0 = \{a,b,c,d,e,f\}$, let $\mu_1 = \mu_1'$ be the identity on
$\{a,b,c,d\}$ (retention $4/6$), let $\mu_2$ be the identity on $\{a,b\}$ and
$\mu_2'$ the identity on $\{c,d\}$, each of retention $2/4$. Both chains
survive $2$ of $6$. Now instead let $\mu_1'$ be the identity on
$\{c,d,e,f\}$ --- still retention $4/6$ --- with $\mu_2'$ the identity on
$\{c,d\}$ of retention $2/4$; the primed chain still survives $\{c,d\}$. Finally
take $\mu_2''$ to be the identity on $\{e,f\}$, retention $2/4$ on
$S_1' = \{c,d,e,f\}$: the chain $\mu_1'$ then $\mu_2''$ survives $\{e,f\}$,
whereas the chain $\mu_1$ then $\mu_2''$ has $S_1 = \{a,b,c,d\}$ meeting
$\operatorname{dom}\mu_2'' = \{e,f\}$ in the empty set, surviving fraction $0$.
The two chains have identical pairwise retentions $(4/6, 2/4)$ and end-to-end
surviving fractions $2/6$ and $0$.
\end{proof}

\subsection{Order does not recover loss}

A natural hope is that a plan losing too much to translation can be repaired by
reordering: translate later, filter earlier, cross the boundary at a different
point. The hope is false in a precise sense.

\begin{theorem}[Reordering does not recover translation loss]
\label{thm:reorder}
Let $\Plan$ contain a translation step $\mu$ and a source step $\sigma$ whose
denotation is a selection $\pi$, and let $\Plan'$ be the plan obtained by
exchanging their order where both orders are well formed. Then
\[
\mu(\pi(S)) \subseteq \mu(S) \cap \pi'(\mu(S))
\quad\text{and}\quad
\pi'(\mu(S)) \subseteq \mu(S),
\]
where $\pi'$ is the image selection. In particular neither order retains an
element of $S$ that lies outside $\operatorname{dom}\mu$, so the two orders
have the same upper bound $r_\mu(S)$ on surviving fraction, and reordering
cannot increase it.
\end{theorem}

\begin{proof}
Both containments are immediate from monotonicity of images and selections.
For the substantive claim, let $u \in S \setminus \operatorname{dom}\mu$. In
the order ``select then translate'', $u$ either fails $\pi$ and is dropped, or
passes $\pi$ and then has no image under $\mu$; either way it contributes
nothing to the output. In the order ``translate then select'', $u$ has no image
under $\mu$ and so contributes nothing before $\pi'$ is applied. Hence in both
orders the output is contained in $\mu(S \cap \operatorname{dom}\mu)$, whose
preimage in $S$ has size at most $r_\mu(S)|S|$. Since this bound is independent
of the order, no exchange increases it.
\end{proof}

\begin{remark}[What reordering does change]
\label{rem:reorder-cost}
\Cref{thm:reorder} says reordering cannot recover \emph{coverage}. It says
nothing about \emph{cost}, which it changes substantially: translating after a
restrictive selection sends fewer identifiers across the boundary and so costs
fewer requests, whereas translating first may allow a batched request where the
other order does not. The planner is therefore free to reorder for cost, and
\cref{thm:allocation} governs that freedom --- but a plan that reorders hoping
to retain more is mistaken, and the language does not offer the operation as a
coverage remedy.
\end{remark}

\subsection{Routes may disagree without either being wrong}

\begin{definition}[Route]
\label{def:route}
A \emph{route} from namespace $n$ to namespace $n'$ is a path in the directed
graph whose vertices are namespaces and whose edges are available translation
maps. Two routes are \emph{parallel} if they share their endpoints.
\end{definition}

\begin{theorem}[Route divergence measures resolved extent]
\label{thm:route-extent}
Let $R = \mu_k \circ \cdots \circ \mu_1$ and $R' = \mu'_l \circ \cdots \circ \mu'_1$
be parallel routes from $n$ to $n'$, each composed of maps that are
\emph{sound} --- every pair they assert is a true correspondence --- and let
$S \subseteq n$. Then:
\begin{enumerate}[label=(\alph*),leftmargin=2.4em]
\item $R(S)$ and $R'(S)$ cannot contradict: every element of each is a true
correspondent of some element of $S$, so neither set contains an element the
other must exclude;
\item the symmetric difference $R(S) \bigtriangleup R'(S)$ consists exactly of
correspondences resolved by one route and unresolved by the other, so
$|R(S) \bigtriangleup R'(S)|$ is a lower bound on the number of correspondences
that at least one route fails to resolve;
\item if every map on both routes is total on the realised sets then
$R(S) = R'(S)$, so divergence implies partiality somewhere;
\item consequently $R(S) \cup R'(S)$ is sound and has surviving fraction at
least $\max(\rho_R(S), \rho_{R'}(S))$.
\end{enumerate}
\end{theorem}

\begin{proof}
(a) By soundness of each constituent map and transitivity of true
correspondence, every $v \in R(S)$ stands in true correspondence to some
$u \in S$, and likewise for $R'$. A set of true correspondents cannot
contradict another such set: contradiction would require one route to assert a
pair the other denies, and neither route denies anything --- both are
existential.

(b) Let $v \in R(S) \setminus R'(S)$. By (a), $v$ is a true correspondent of
some $u \in S$, and $R'$ does not produce it, so $R'$ fails to resolve the
correspondence $(u,v)$. Symmetrically for the other difference. Each element of
the symmetric difference therefore witnesses at least one unresolved
correspondence on one route, and distinct elements witness distinct
correspondences, giving the bound.

(c) If every map is total on the realised sets, each route computes the full
composite correspondence relation restricted to $S$, and the composite relation
is determined by the underlying true correspondence, which is route-independent.
So the two images coincide.

(d) The union of two sound sets is sound, and contains each, so its surviving
fraction is at least each of theirs.
\end{proof}

\begin{remark}[The union is not free]
Part (d) suggests always taking the union of parallel routes. The cost is that
the union costs the sum of the routes' requests, and by
\cref{thm:allocation} that budget comes from somewhere. The language exposes
\texttt{union} over routes precisely so that the trade is written down rather
than assumed, and \cref{ex:route-union} shows the syntax.
\end{remark}

\begin{example}[Two routes, written down]
\label{ex:route-union}
\begin{lstlisting}[caption={Parallel routes made explicit. The plan asserts nothing about which is better; it measures the divergence and reports it.},label=lst:routes]
let direct   = map compounds via chebi_to_kegg
let indirect = map compounds via chebi_to_inchikey
               then via inchikey_to_kegg

let both     = union direct indirect
assert soundness(direct) and soundness(indirect)
emit divergence(direct, indirect) as resolved_extent
\end{lstlisting}
By \cref{thm:route-extent}(b) the emitted \texttt{resolved\_extent} is a lower
bound on the correspondences at least one route misses. It is not an error
count and must not be reported as one: both routes may be individually correct
and the divergence is then a statement about coverage, not about accuracy.
\end{example}

\subsection{The retention check}

\begin{definition}[Retention check]
\label{def:retention-check}
A translation step annotated \texttt{expect partial} $\epsilon$ computes
$r_\mu(S)$ on its actual input and, if $r_\mu(S) < \epsilon$, returns
$\vstarve$ to its successors with a diagnosis naming $\mu$, $\epsilon$ and the
observed $r_\mu(S)$. Otherwise it returns $\vans$ with the image and records
both $r_\mu(S)$ and $a_\mu(S)$ as provenance attributes.
\end{definition}

\begin{proposition}[The check converts a silent decimation into a located verdict]
\label{prop:check-locates}
Let $\Plan$ contain a translation step $\Step_t$ with retention below its
declared $\epsilon$, feeding a step $\Step_d$ that would otherwise report
$\vempty$. Under \cref{def:retention-check}, $\Step_d$ reports $\vstarve$ and
the blame chain of \cref{prop:blame} terminates at $\Step_t$.
\end{proposition}

\begin{proof}
By \cref{def:retention-check}, $\Step_t$ returns a non-$\vans$ verdict. By (R1)
of \cref{def:verdicts}, $\Step_d$ reports $\vstarve$ naming $\Step_t$. By
\cref{prop:blame} the chain terminates at a non-$\vstarve$ step; since
$\Step_t$'s verdict was assigned by \cref{def:retention-check} rather than by
(R1), and $\Step_t$ names no predecessor, the chain terminates there.
\end{proof}

This is the resolution of defect (D3) in its commonest concrete form. Without
the check, the analyst sees an empty pathway result and investigates the
pathway database. With it, the executor reports that the translation two steps
earlier retained $0.31$ against a declared $0.6$, and names the map.

%==============================================================================
\part{Execution}
%==============================================================================

\section{Budget allocation across heterogeneous sources}
\label{sec:allocation}

\subsection{Why allocation is forced}

The sources of \cref{obs:kind} differ in cost by orders of magnitude and in
kind of limit. A locally loaded artefact is effectively free and bounded by
memory; a SPARQL endpoint is bounded by a per-request timeout and an
unadvertised fair-use policy; a flat-file REST interface is bounded by an
explicit request rate. A plan that ignores this either underuses the cheap
sources or exceeds the expensive ones, and \cref{lst:script} does both. Since
the total request budget is finite and the sources' yields are different
functions of the effort spent on them, allocation is an optimisation problem
and not a scheduling convention.

\begin{definition}[Yield]
\label{def:yield}
The \emph{yield} of step $i$ under effort $e_i \geq 0$ is
$\gamma_i(e_i)$, where $\gamma_i : \Rset_{\geq 0} \to \Rset_{\geq 0}$ is
non-decreasing, continuously differentiable, strictly concave, and satisfies
$\gamma_i(0) = 0$. Effort is measured in requests.
\end{definition}

Strict concavity is the substantive assumption and it is the empirically
plausible one: the first requests to a source retrieve the common cases and
later requests retrieve progressively rarer ones. It fails for a step whose
result is all-or-nothing --- a single lookup that either succeeds or does not
--- and \cref{rem:concavity-fails} records the consequence.

\begin{definition}[Allocation problem]
\label{def:alloc}
Given a plan with $m$ steps, total budget $\Bud$, and per-step minimum costs
$c_i > 0$, the allocation problem is
\[
\max_{e \in \Rset^m_{\geq 0}} \sum_{i=1}^{m} \gamma_i(e_i)
\quad\text{subject to}\quad
\sum_{i=1}^{m} e_i \leq \Bud .
\]
\end{definition}

\subsection{A single shadow price}

\begin{theorem}[Water-filling allocation]
\label{thm:allocation}
The allocation problem of \cref{def:alloc} has a unique optimum $e^\ast$, and
there is a scalar $p^\ast \geq 0$ --- the shadow price of a request --- such
that for every $i$,
\[
e_i^\ast > 0 \implies \gamma_i'(e_i^\ast) = p^\ast,
\qquad
e_i^\ast = 0 \implies \gamma_i'(0) \leq p^\ast ,
\]
and $\sum_i e_i^\ast = \Bud$ whenever some $\gamma_i'(0) > 0$. Consequently
$e_i^\ast = (\gamma_i')^{-1}(p^\ast)$ on the support, and $p^\ast$ is the unique
root of $\sum_i \max\bigl(0, (\gamma_i')^{-1}(p)\bigr) = \Bud$.
\end{theorem}

\begin{proof}
The objective is a sum of strictly concave functions, hence strictly concave;
the feasible set is a compact convex simplex-like region. A strictly concave
function on a compact convex set attains a unique maximum, giving existence and
uniqueness.

Form the Lagrangian
$L(e,p,\lambda) = \sum_i \gamma_i(e_i) - p(\sum_i e_i - \Bud) + \sum_i \lambda_i e_i$
with $p \geq 0$ the multiplier of the budget constraint and $\lambda_i \geq 0$
those of the non-negativity constraints. Slater's condition holds --- the point
$e_i = \Bud/(2m)$ is strictly feasible --- so the KKT conditions are necessary
and sufficient. Stationarity gives
$\gamma_i'(e_i^\ast) - p^\ast + \lambda_i^\ast = 0$ for each $i$. Complementary
slackness gives $\lambda_i^\ast e_i^\ast = 0$. If $e_i^\ast > 0$ then
$\lambda_i^\ast = 0$ and $\gamma_i'(e_i^\ast) = p^\ast$. If $e_i^\ast = 0$ then
$\lambda_i^\ast \geq 0$ and $\gamma_i'(0) = p^\ast - \lambda_i^\ast \leq p^\ast$.

If some $\gamma_i'(0) > 0$ then the budget constraint is active at the optimum:
were $\sum_i e_i^\ast < \Bud$, adding $\delta > 0$ to that $e_i$ would remain
feasible and, by continuity of $\gamma_i'$ and $\gamma_i'(0)>0$, would strictly
increase the objective, contradicting optimality. Hence $\sum_i e_i^\ast = \Bud$.

Strict concavity makes each $\gamma_i'$ strictly decreasing and therefore
invertible on its range, so $e_i^\ast = (\gamma_i')^{-1}(p^\ast)$ where
positive, and $\sum_i \max(0,(\gamma_i')^{-1}(p))$ is continuous and strictly
decreasing in $p$ on the region where it is positive, hence has a unique root
at $\Bud$.
\end{proof}

\begin{corollary}[Nothing is dropped by a rule]
\label{cor:no-dropping-rule}
A step receives zero effort at the optimum if and only if
$\gamma_i'(0) \leq p^\ast$, that is, if and only if its marginal yield at the
first request is below the price. No step is excluded by a selection heuristic,
a priority ordering, or a source blacklist.
\end{corollary}

\begin{proof}
Immediate from the two implications of \cref{thm:allocation}, which are
exhaustive and mutually exclusive given strict concavity.
\end{proof}

\Cref{cor:no-dropping-rule} matters because the alternative --- a hand-written
priority list saying which sources to try first --- is what \cref{lst:script}
does implicitly through its statement order, and it is unfalsifiable: one
cannot ask of a priority list whether it is right. One can ask of $p^\ast$
whether the yields that determined it were estimated correctly, which is a
question with an answer.

\begin{proposition}[Sufficient total budget is not sufficient]
\label{prop:necessary-not-sufficient}
Let $c_i$ be the minimum cost at which step $i$ can return a non-empty payload.
Then $\Bud \geq \sum_i c_i$ is necessary for a plan to complete with all steps
$\vans$, and is not sufficient.
\end{proposition}

\begin{proof}
Necessity: each step must be issued at cost at least $c_i$ to return a
non-empty payload, and the costs are additive over steps and drawn from one
budget, so any completing execution spends at least $\sum_i c_i$.

Insufficiency: take $m=2$ with $c_1 = c_2 = 1$ and $\Bud = 2$. Let step $2$
consume the output of step $1$, and let $\cost_{\Src_2}$ be increasing in input
cardinality with $\cost_{\Src_2}(k) = k$. If step $1$ at cost $1$ returns a
payload of cardinality $3$, then step $2$ requires cost $3$, but only $1$
remains. The plan halts with step $2$ reporting $\vref$ by rule (R3) of
\cref{def:verdicts}, despite $\Bud \geq c_1 + c_2$. The obstruction is that
$c_2$ is a minimum over inputs and the realised input is not the minimising
one, and no static quantity bounds the realised input without executing step
$1$.
\end{proof}

\begin{remark}[The consequence for planning]
\label{rem:replanning}
\Cref{prop:necessary-not-sufficient} shows that a purely static allocation can
be infeasible at run time through no error of the allocator. Two responses are
available. The first is to re-solve \cref{thm:allocation} after each step with
the remaining budget and the realised cardinalities, which is correct and costs
one convex solve per step. The second is to allocate statically from upper
bounds on cardinality, which is cheap and wastes budget in proportion to how
loose the bounds are. The prototype takes the first, and \cref{sec:limits}
records that we have not established a regret bound for it against the
clairvoyant allocation --- the natural comparison, which we do not make.
\end{remark}

\begin{remark}[Where concavity fails]
\label{rem:concavity-fails}
A \textsf{lookup} step against a REST interface has an all-or-nothing yield:
one request retrieves the record and further requests retrieve nothing. Its
yield is a step function, neither strictly concave nor differentiable, and
\cref{thm:allocation} does not apply to it. The prototype handles such steps by
removing them from the optimisation, charging their fixed cost to the budget
first, and solving \cref{def:alloc} over the remainder --- which is correct
because their effort is not a decision variable, and which means the shadow
price governs only the steps that have a real allocation choice.
\end{remark}

\section{Structural interpolation}
\label{sec:interp}

\subsection{An instance of the defect}

The following is a documented observation and we restate it in full because the
argument of this section depends on the details rather than on the headline.

\begin{observation}[Two spellings, two answers]
\label{obs:2-397}
On 10 August 2026, two requests were issued to a public SPARQL endpoint of a
reaction knowledge base, with a JSON results media type. The two requests were
byte-identical except for one line. The first, using an abbreviated object
list, returned a count of $2$. The second, using two separate triple patterns,
returned a count of $397$. The ratio is $198.5$.
\end{observation}

\begin{lstlisting}[caption={Form A. Returns 2. The abbreviation is on line 6.},label=lst:comma,language=,keywordstyle=]
PREFIX rh:    <http://rdf.rhea-db.org/>
PREFIX rdfs:  <http://www.w3.org/2000/01/rdf-schema#>
PREFIX chebi: <http://purl.obolibrary.org/obo/CHEBI_>
SELECT (COUNT(DISTINCT ?r) AS ?n) WHERE {
  ?r rdfs:subClassOf rh:Reaction ; rh:status rh:Approved .
  ?r rh:side/rh:contains/rh:compound/rh:chebi ?a , ?o .
  VALUES ?aa { chebi:35238 chebi:37022 }
  VALUES ?ox { chebi:35179 chebi:36147 chebi:133294 }
  ?a rdfs:subClassOf* ?aa .
  ?o rdfs:subClassOf* ?ox .
}
\end{lstlisting}

\begin{lstlisting}[caption={Form B. Returns 397. Differs from \cref{lst:comma} only in that line 6 is written as two patterns.},label=lst:longhand,language=,keywordstyle=]
  ?r rh:side/rh:contains/rh:compound/rh:chebi ?a .
  ?r rh:side/rh:contains/rh:compound/rh:chebi ?o .
\end{lstlisting}

The two forms are semantically equivalent under the specification. An object
list is syntax: the grammar admits a comma-separated object list in a
predicate-object list, and the abbreviation is expanded during translation to
the abstract query --- before evaluation, and without a special case for
property paths \citep{Harris2013}. A conforming implementation must therefore
return the same count for both. The difference is not a specification
ambiguity.

\begin{remark}[The control is clean and the confound is named]
\label{rem:control}
Both forms were evaluated locally against two independent engines
\citep{RDFLib,Pyoxigraph} over a hand-checkable dataset, and both engines
returned the hand-computed answer for both spellings, so the equivalence is not
merely our reading of the grammar. Separately, the endpoint's loaded ontology
carried $204{,}585$ classes against $237{,}672$ in the corresponding published
download --- a difference of $33{,}087$, or $13.9\%$ --- so any comparison of
counts \emph{across} the two artefacts is confounded by snapshot skew. That
confound does not touch \cref{obs:2-397}, in which both counts come from the
same endpoint in the same session.
\end{remark}

\begin{remark}[We do not diagnose the service]
\label{rem:no-probing}
The cause lies in the endpoint's translation of the abbreviated form and we do
not attempt to determine it further. Establishing it would require probing a
third party's public service with requests designed to elicit internal
behaviour, which is not ours to do. Nothing in this section should be read as a
recommendation to do so, and the design consequence below does not depend on
knowing the cause.
\end{remark}

\subsection{The design consequence}

\Cref{obs:2-397} is an instance of a family. In each member, a plan author
writes a request that is correct under the target language's specification, a
deployment interprets one spelling differently from another, and the difference
is invisible because it is a difference between two strings the author regards
as the same string. Textual interpolation puts the author in exactly this
position, because the string that is sent is assembled at run time and never
inspected.

\begin{construction}[Longhand-only lowering]
\label{cons:longhand}
The lowering $\Low_\Src$ emits, for each abstract request, a canonical concrete
form fixed by the adapter, in which (i) every predicate-object list is expanded
to separate patterns, (ii) every bound result set enters through the target
language's own value-binding construct rather than by concatenation into the
pattern text, and (iii) no construct is emitted whose expansion interacts with
another construct in the sense of \cref{rem:independence-fails}.
\end{construction}

\begin{theorem}[Structural interpolation eliminates the defect family]
\label{thm:interpolation}
Let $\mathcal{D}$ be the family of defects in which two concrete requests with
the same denotation under the target specification are assigned different
denotations by a deployment, and whose difference lies in the spelling of a
construct that \cref{cons:longhand} canonicalises. Then no execution of a plan
under \cref{cons:longhand} exhibits a member of $\mathcal{D}$.
\end{theorem}

\begin{proof}
Fix a plan and an execution. Every concrete request issued is produced by
$\Low_\Src$, since by \cref{prin:leaves} no other source of concrete requests
exists in the language. By \cref{cons:longhand}, $\Low_\Src$ emits a single
canonical spelling for each abstract request: the expanded, longhand form with
bindings supplied structurally. Membership in $\mathcal{D}$ requires two
requests differing in a canonicalised spelling. But for any two abstract
requests with the same denotation, the emitted concrete forms are determined by
the canonical procedure and therefore agree wherever the canonicalisation
applies; and no non-canonical spelling is ever emitted, so no pair of emitted
requests differs in one. Hence no member of $\mathcal{D}$ is exhibited.
\end{proof}

\begin{remark}[Precisely what is and is not eliminated]
\label{rem:interp-scope}
\Cref{thm:interpolation} is a construction-by-exclusion result and its strength
is exactly the coverage of \cref{cons:longhand}. It eliminates the family
$\mathcal{D}$ --- differences attributable to a spelling the lowering
canonicalises. It does not eliminate: deployments that misinterpret the
canonical form itself; capability over-declaration
(\cref{rem:honesty-assumption}), which is a different defect with the same
symptom; or divergence between a service and a specification in constructs the
canonicalisation does not touch. The honest statement is that the language
removes the author's ability to make one class of mistake, by removing the
author's access to the string, and that this is a smaller claim than
correctness.
\end{remark}

\begin{proposition}[The elimination is not available to textual interpolation]
\label{prop:textual-cannot}
Let $I$ be an interface in which the caller supplies the concrete request as a
string. Then for any canonicalisation procedure $C$, there is a caller of $I$
whose issued request is not in the image of $C$.
\end{proposition}

\begin{proof}
The caller supplies an arbitrary string and $I$ transmits it. Since $C$'s image
is a proper subset of the concrete language --- proper because $C$
canonicalises, so at least two distinct strings map to one, and the
non-canonical one is outside the image --- a caller supplying that
non-canonical string issues a request outside the image. No property of $I$
prevents this, because $I$'s contract is to transmit what it is given.
\end{proof}

\Cref{prop:textual-cannot} is why \cref{prin:leaves} is a principle of the
language rather than a recommended practice: a recommendation is exactly the
thing \cref{prop:textual-cannot} says is unenforceable.

\section{Re-execution}
\label{sec:idempotence}

A plan is often run more than once --- to refresh a result, to recover from a
budget exhaustion, to check a suspicious answer. Whether re-execution is safe
is not obvious, because a plan mixes sources with different volatility and a
retention check compares against a threshold.

\begin{definition}[Stability]
\label{def:stable}
A source $\Src$ is \emph{stable over an interval} if the dataset it holds is
unchanged throughout. A translation map is stable if its relation is unchanged.
\end{definition}

\begin{theorem}[Characterisation of safe re-execution]
\label{thm:idempotent}
Let $\Plan$ be well-capability with total budget $\Bud$, executed twice over an
interval throughout which every source and every translation map it uses is
stable. Then the two executions assign the same verdict and the same payload to
every step if and only if both executions allocate every step an effort at
least its realised cost. If some step is allocated less in one execution than
its realised cost, the two executions may differ, and the difference is
confined to that step and its successors in $G(\Plan)$.
\end{theorem}

\begin{proof}
($\Leftarrow$) Suppose both executions allocate every step at least its
realised cost. We induct along the sequence order. For $\Step_1$,
$\beta_1 = ()$, so (R1) does not fire; $\Req(\rho_1) \subseteq \Capset(\Src_1)$ by
well-capability, so (R2) does not fire; the allocation covers the realised
cost, so neither (R3) nor (R4) fires; the request is issued, and by stability
the dataset is the same in both executions, so by faithfulness
(\cref{thm:lowering}) the extracted result equals
$\llbracket \rho_1 \rrbracket_D$ in both. Hence the verdict and payload agree.
For the inductive step, assume all predecessors of $\Step_i$ agree across the
two executions. Then the antecedent of (R1) has the same truth value in both,
and if it is false the same argument as the base case applies, using stability
of any translation map involved and the fact that a retention check compares a
fixed $r_\mu(S)$ against a fixed $\epsilon$, so it too agrees.

($\Rightarrow$) Suppose some step is allocated less than its realised cost in
some execution. Then in that execution rule (R3) or (R4) fires at that step,
yielding $\vref$ or $\vtime$; if the other execution allocates at least the
realised cost, that step yields $\vans$ or $\vempty$ there. The verdicts
differ, contradicting agreement. Hence agreement forces sufficient allocation
everywhere.

Confinement: by \cref{def:verdicts}, a step's verdict depends only on its own
source, capabilities, budget and the verdicts and payloads of its predecessors.
If all predecessors of $\Step_j$ agree and $\Step_j$ is not the differing step,
the same argument as above gives agreement at $\Step_j$. So the difference
propagates only along edges of $G(\Plan)$ out of the differing step.
\end{proof}

\begin{corollary}[Budget exhaustion is recoverable, data change is not]
\label{cor:rerun}
Re-running a plan that halted with $\vref$ or $\vtime$, under a larger budget
and over a stable interval, reproduces every verdict the first execution
assigned before the halt and may extend the plan further. Re-running across a
change in any source's dataset carries no such guarantee, and the language
records source snapshot identifiers in provenance so that the two cases are
distinguishable after the fact.
\end{corollary}

\begin{proof}
The first claim is \cref{thm:idempotent}($\Leftarrow$) applied to the prefix of
steps allocated sufficiently in both executions, which by hypothesis includes
every step completed before the halt. The second follows because stability is a
hypothesis of \cref{thm:idempotent} and the theorem's conclusion is not
asserted without it; the provenance recording is a design decision, not a
consequence.
\end{proof}

%==============================================================================
\part{Realisation}
%==============================================================================

\section{A prototype}
\label{sec:proto}

\subsection{What the prototype is, and what it is not}

The prototype is a compiler and executor for \HFQ{} written in Python. It parses
plan text into the intermediate representation of \cref{def:plan}, performs the
static check of \cref{thm:static}, solves the allocation of
\cref{thm:allocation}, executes the plan against source adapters, assigns
verdicts by \cref{def:verdicts}, and writes the result --- verdicts, payloads,
retention figures, allocation, and provenance --- as a JSON document.

It is not a benchmark. Every adapter in the prototype resolves against a local
fixture or a local engine. No request leaves the machine. This is a deliberate
constraint and not a temporary limitation of the implementation: the properties
under test are properties of the compiler and the verdict assignment, and a
live service can neither confirm nor refute them while introducing a dependency
on a third party's availability, policy, and current snapshot. What a live
service would test --- whether real cross-reference tables have the retention
figures one hopes --- is a question this paper does not answer and does not
claim to (\cref{sec:limits}).

\subsection{Pipeline}

\begin{enumerate}[label=(\arabic*),leftmargin=2.4em]
\item \textbf{Parse.} Plan text to a list of steps
$(x_i, \Src_i, \rho_i, \beta_i, b_i)$, together with the plan budget. The
grammar is small: a plan is a budget declaration followed by \texttt{let}
bindings and a terminal \texttt{emit}.
\item \textbf{Resolve.} Each source name is resolved against a registry of
adapters, each of which declares $\Capset(\Src)$ as an explicit set of feature
symbols. The declaration is data, written by the adapter author, and is the
locus of the honesty assumption (\cref{rem:honesty-assumption}).
\item \textbf{Check.} Compute $\Req(\rho_i)$ for each step by structural
recursion over the abstract request, and test
$\Req(\rho_i) \subseteq \Capset(\Src_i)$. On failure the executor halts before
issuing any request and emits a refusal document naming the missing features
and the step --- \cref{cor:refuse-before-contact} made operational.
\item \textbf{Allocate.} Solve \cref{def:alloc} by bisection on the shadow price
$p$, using $\gamma_i(e) = w_i \log(1 + e)$ as the default yield, whose
derivative $w_i/(1+e)$ inverts in closed form to
$e = w_i/p - 1$. Steps with step-function yields are charged first
(\cref{rem:concavity-fails}).
\item \textbf{Execute.} In sequence order, apply (R1)--(R6). A step reaching
(R4) calls its adapter's lowering, which emits the canonical concrete request of
\cref{cons:longhand}; the adapter evaluates it against its fixture and returns a
result set.
\item \textbf{Emit.} Serialise to JSON.
\end{enumerate}

\subsection{Adapters}

Four adapter kinds are implemented, chosen to exercise the four source kinds of
\cref{obs:kind} rather than to be complete.

\begin{itemize}[leftmargin=2.4em]
\item A \emph{graph-pattern adapter} over a local RDF store
\citep{RDFLib,Pyoxigraph}, declaring
$\{\textsf{pattern}, \textsf{path}, \textsf{bind}, \textsf{filter}, \textsf{agg}\}$.
Its lowering emits longhand triple patterns and supplies bound sets through the
value-binding construct, never by concatenation.
\item A \emph{lookup adapter} over a local key--value fixture, declaring
$\{\textsf{lookup}, \textsf{link}\}$ and \emph{not} \textsf{pattern}. It is the
prototype's stand-in for a flat-file REST interface, and its restricted
declaration is what makes \cref{thm:static} fire on plans that ask it for a
join.
\item An \emph{ontology adapter} over a local class hierarchy, declaring
$\{\textsf{path}, \textsf{lookup}\}$ --- transitive closure over subsumption but
no filtering and no aggregation.
\item A \emph{map adapter} implementing translation steps
(\cref{def:map}), which is the only adapter that reports retention and
amplification.
\end{itemize}

\begin{remark}[The declarations are the experiment]
\label{rem:decl-experiment}
Because the static check is decided entirely by the declarations, and the
declarations are chosen by us, the prototype cannot demonstrate that the check
is \emph{right} about any real service. What it demonstrates is that the check
is \emph{decidable and cheap}, that refusal precedes contact, and that the
verdict assigned to a rejected plan names the missing feature rather than
reporting an empty result. Those are the claims of \cref{thm:static} and
\cref{cor:refuse-before-contact}, and they are properties of the compiler.
\end{remark}

\subsection{The output document}

The JSON document produced by an execution has one entry per step, and the
schema is fixed by the theory rather than by convenience:

\begin{lstlisting}[caption={Shape of a step entry in the emitted JSON. The \texttt{blocker} field is absent exactly when the verdict is \vempty{} or \vans{}, per \cref{def:blocker}.},label=lst:json,language=,keywordstyle=]
{
  "step": "kids",
  "source": "CHEBI->KEGG",
  "verdict": "starved",
  "blocker": "corpus",
  "diagnosis": {
    "named_predecessor": "ecs",
    "reason": "retention 0.31 below declared expectation 0.6"
  },
  "retention": 0.31,
  "amplification": 1.14,
  "budget": {"allocated": 60, "spent": 60, "shadow_price": 0.0142},
  "payload": null,
  "provenance": {"snapshot": "fixture-v3", "lowered_form": "canonical"}
}
\end{lstlisting}

Three features of \cref{lst:json} are consequences of results above rather than
design taste. The \texttt{blocker} field is partial, because
\cref{def:blocker} is partial and $\vempty$ has no blocker. The
\texttt{payload} is \texttt{null} for every non-$\vans$ verdict, because
\cref{prin:verdict} says a payload accompanies only $\vans$ --- and an
implementation that returned a partial payload alongside a failure verdict
would reintroduce exactly the ambiguity \cref{cor:onebit} identifies.
\texttt{retention} and \texttt{amplification} are recorded separately, because
by \cref{prop:cardinality-uninformative} their product --- which is all the
output cardinality reveals --- determines neither.

\section{Validation design}
\label{sec:validation}

We enumerate the checks the prototype performs, in the interest of making it
clear which claims of this paper are exercised and which are not. Each is
executed against local fixtures.

\begin{enumerate}[label=(V\arabic*),leftmargin=3em]
\item \textbf{Static check decides without contact.} For a family of plans,
half well-capability and half not, verify that the ill-capability plans halt
with a refusal document, that the refusal names the correct missing features,
and that the adapter request counter is zero.
(\cref{thm:static}, \cref{cor:refuse-before-contact}.)
\item \textbf{Check cost is linear.} Measure the check's operation count
against $m|\Feat|$ over plans of growing length. (\cref{thm:static}(a).)
\item \textbf{Six verdicts are realised.} Construct the six configurations of
the proof of \cref{thm:six} --- each differing from a common baseline in one
respect --- and verify that the executor assigns six distinct verdicts.
(\cref{thm:six}.)
\item \textbf{The one-bit interface conflates four.} Run the same six
configurations through a deliberately impoverished adapter that returns only a
result set, and verify that four of the six are indistinguishable in its
output. (\cref{cor:onebit}.)
\item \textbf{Starvation is unreachable when $m=1$.} Over all single-step
configurations in the fixture family, verify that $\vstarve$ is never assigned
and $\bcorp$ never emitted; then verify that the predecessor perturbation of
\cref{prop:starve-reachable}(b) produces it in a two-step plan.
\item \textbf{Blame terminates.} For every starved execution, follow the
diagnosis chain and verify it terminates at a non-$\vstarve$ step within $m$
hops. (\cref{prop:blame}.)
\item \textbf{Retention factorises.} Over random partial maps, verify the
multiplicative identity of \cref{thm:retention}(a) exactly, and verify that the
measured surviving fraction lies within the bounds of (b) and (c). Record how
often each bound is attained and how wide the gap is
(\cref{rem:bounds-gap}).
\item \textbf{Cardinality is uninformative.} Realise the two maps of
\cref{prop:cardinality-uninformative} and verify equal output cardinality with
retentions differing by the predicted factor.
\item \textbf{Pairwise coverage does not determine chain coverage.} Realise the
two families of \cref{prop:pairwise-insufficient} and verify identical pairwise
retentions with differing end-to-end surviving fractions.
\item \textbf{Reordering does not recover coverage.} For plans admitting both
orders, verify equal surviving fractions and unequal request counts.
(\cref{thm:reorder}, \cref{rem:reorder-cost}.)
\item \textbf{Routes diverge without contradiction.} Construct parallel sound
routes with partial maps, verify that neither contains an element contradicting
the other, that the symmetric difference is non-empty, and that making every
map total collapses the divergence to zero.
(\cref{thm:route-extent}.)
\item \textbf{Allocation satisfies KKT.} Solve \cref{def:alloc} numerically and
verify that the marginal yields of all supported steps agree to tolerance, that
unsupported steps satisfy $\gamma_i'(0) \leq p^\ast$, and that the budget binds.
(\cref{thm:allocation}.)
\item \textbf{Sufficient budget is not sufficient.} Realise the two-step
counterexample of \cref{prop:necessary-not-sufficient} and verify a $\vref$
verdict despite $\Bud \geq \sum c_i$.
\item \textbf{Lowering is canonical.} For every abstract request in the fixture
family, verify that the emitted concrete form contains no predicate-object list
and no interpolated identifier outside a value-binding construct, and that two
abstract requests with the same denotation lower to the same string.
(\cref{cons:longhand}, \cref{thm:interpolation}.)
\item \textbf{Re-execution is stable.} Execute each plan twice over an unchanged
fixture with sufficient budget and verify identical verdicts and payloads; then
re-execute with a budget below one step's realised cost and verify that the
difference is confined to that step and its successors.
(\cref{thm:idempotent}.)
\end{enumerate}

We record what these do \emph{not} establish, since the enumeration invites the
opposite reading. (V1)--(V15) test the compiler and the executor against the
theory. They do not test the theory against the world. No check above would
detect a dishonest capability declaration, an inaccurate cross-reference table,
or a deployment that misinterprets the canonical form --- and the first and
third of those are precisely the failure modes \cref{rem:honesty-assumption}
and \cref{rem:interp-scope} identify as beyond the reach of the construction.

\section{Limitations}
\label{sec:limits}

\paragraph{Honesty is assumed, never verified.} The whole of \cref{thm:static}
and the greater part of \cref{sec:verdicts} rest on the assumption that a
source's declared capability set is honest (\cref{def:honest}). Under-declaration
costs expressiveness and is safe; over-declaration is unsound and, worse,
invisible --- an over-declared source accepts a request it will silently degrade,
and the degraded answer arrives as $\vans$ with a payload. We do not know how to
verify a declaration without probing the service, and probing a third party's
public service to characterise its internals is not something we undertake.
This is the single largest gap in the development and we do not minimise it.

\paragraph{Yields are posited, not estimated.} \Cref{thm:allocation} assumes
strictly concave differentiable yield functions and the prototype supplies
$w_i\log(1+e)$ by default. We have not estimated a yield function for any real
source, and the weights $w_i$ are configuration. An allocation optimal for the
wrong yields is not optimal. Furthermore \cref{rem:concavity-fails} excludes
all-or-nothing steps from the optimisation entirely, so the shadow price
governs a subset of the plan.

\paragraph{No regret bound for replanning.} \Cref{rem:replanning} adopts
re-solving after each step, and we have not bounded its total yield against the
clairvoyant allocation that knew the realised cardinalities in advance. This is
the natural theoretical question raised by
\cref{prop:necessary-not-sufficient} and we do not answer it.

\paragraph{Retention laws are unmeasured.} \Cref{thm:retention} is arithmetic
and holds of any partial maps. Whether real cross-reference tables have
retentions making \cref{prop:starve-common} bite, and whether the losses of
successive maps are correlated or independent, is an empirical question we
raise and do not settle. We note only that the independence hypothesis is
false in the direction that makes the bound conservative.

\paragraph{Namespace disjointness is a modelling decision.} \Cref{rem:disjoint}
records that we take namespaces to be pairwise disjoint. Real identifier spaces
overlap and are re-used, and a shared identifier appearing in two namespaces
would be conflated by the model. We do not know how badly.

\paragraph{Wall-clock time is not modelled.} All costs count requests
(\cref{sec:intro}, conventions). A plan whose requests are few and slow and one
whose requests are many and fast are indistinguishable to the allocator. The
verdict $\vtime$ mentions a clock but the budget does not, and reconciling the
two would require a second resource dimension the model does not carry.

\paragraph{The static check is not a completeness result.} \Cref{thm:static}
says a well-capability plan compiles. It does not say an ill-capability plan is
unanswerable: a source might answer a request outside its declared capabilities
through a route the adapter does not expose. Refusal is conservative, and
conservatism is the safe direction, but it is not optimality.

\paragraph{Structural interpolation is a smaller claim than correctness.}
\Cref{rem:interp-scope} states this in place and we repeat it here because the
result is the one most liable to be over-read. \Cref{thm:interpolation}
eliminates a family of defects by removing the author's access to the request
string. It does not make the request correct, and a deployment that misreads the
canonical form is untouched by it.

\paragraph{Nothing has been run against a live service.} By construction. The
consequence is that every quantitative figure in the prototype's output is a
figure about a fixture, and no claim in this paper is supported by measurement
of a public biological database.

\section{Conclusion}
\label{sec:conclusion}

The script of \cref{lst:script} is not bad code. It is the best code available
under a division of labour in which control flow lives in a host language and
data access lives in a query language, with no channel between them. Its five
defects are consequences of that division and not of its author.

We have described a language that moves the division. Control flow and data
access live together in a plan; the plan author never writes a request; requests
are generated by adapters from abstract steps and source declarations. The
consequences follow from that single relocation. Because requests are generated,
their spelling is canonical and a family of interpolation defects cannot occur
(\cref{thm:interpolation}). Because sources declare capabilities, compilability
is decidable before contact and refusal precedes any request
(\cref{thm:static}). Because a step returns a verdict rather than rows, the four
outcomes a row-returning interface conflates become distinguishable
(\cref{cor:onebit}) --- including one, starvation, that has no single-database
analogue and becomes reachable only in federation
(\cref{prop:starve-reachable}). Because translation is a first-class step, its
retention is measured rather than assumed, and the chain laws say what a plan
can and cannot infer about its own coverage
(\cref{thm:retention}, \cref{prop:pairwise-insufficient}) --- including that
reordering will not help (\cref{thm:reorder}) and that two sound routes may
disagree without either being wrong (\cref{thm:route-extent}). Because the
budget is declared, allocation is an optimisation with a single shadow price
rather than a priority list (\cref{thm:allocation}), though a sufficient total
is still not sufficient (\cref{prop:necessary-not-sufficient}).

What the development does not have is a way to check the one thing it most
depends on. A capability declaration is an assertion by an adapter author, and
an over-declaration is unsound and silent. Every static guarantee in this paper
is conditional on that assertion. Whether the condition can be discharged
without probing the services in question --- which is not ours to do --- is the
question we would most like answered and the one we leave open.


\clearpage
\bibliographystyle{plainnat}
\bibliography{references}

\end{document}
